\IfFileExists{stacks-project.cls}{%
\documentclass{stacks-project}
}{%
\documentclass{amsart}
}
\usepackage{amssymb}

% For dealing with references we use the comment environment
\usepackage{verbatim}
\usepackage{fancyvrb}
\usepackage{pmboxdraw}
\newenvironment{reference}{\comment}{\endcomment}
%\newenvironment{reference}{}{}
\newenvironment{slogan}{\comment}{\endcomment}
\newenvironment{history}{\comment}{\endcomment}

% For source figures retained as assets
\usepackage{graphicx}

% For commutative diagrams we use Xy-pic
\usepackage[all]{xy}

% We use 2cell for 2-commutative diagrams.
\xyoption{2cell}
\UseAllTwocells

% We use multicol for the list of chapters between chapters
\usepackage{multicol}

% This is generally recommended for better output
\usepackage{lmodern}
\usepackage[T1]{fontenc}
\usepackage{textalpha}

% For cross-file-references
% We prefer the snapshot tmp/xref/#2.aux,
% which latexmkrc replaces atomically, so
% that documents compiled simultaneously
% never read a half-written .aux file.
% The last argument keeps links to #2.pdf.
\usepackage{xr-hyper}
\newcommand{\maybeexternaldocument}[2]{%
  \IfFileExists{#2.aux}{%
    \externaldocument[#1]{#2}%
  }{%
    \PackageWarning{stack}{Missing #2.aux; external references with prefix #1 unavailable}%
  }%
}
\let\externaldocumentorig\externaldocument
\renewcommand{\externaldocument}[2][]{%
  \IfFileExists{tmp/xref/#2.aux}{%
    \externaldocumentorig[#1]{tmp/xref/#2}[#2.pdf]%
  }{%
  \IfFileExists{#2.aux}{%
    \externaldocumentorig[#1]{#2}%
  }{%
    \PackageWarning{stack}{Missing #2.aux; external references with prefix #1 unavailable}%
  }%
  }%
}

% Package for hypertext links:
\usepackage{hyperref}

% For any local file, say "hello.tex" you want to link to please
% use \externaldocument[hello-]{hello}
\externaldocument[introduction-]{introduction}
\externaldocument[conventions-]{conventions}
\externaldocument[sets-]{sets}
\externaldocument[categories-]{categories}
\externaldocument[topology-]{topology}
\externaldocument[sheaves-]{sheaves}
\externaldocument[sites-]{sites}
\externaldocument[stacks-]{stacks}
\externaldocument[fields-]{fields}
\externaldocument[algebra-]{algebra}
\externaldocument[brauer-]{brauer}
\externaldocument[homology-]{homology}
\externaldocument[derived-]{derived}
\externaldocument[simplicial-]{simplicial}
\externaldocument[more-algebra-]{more-algebra}
\externaldocument[smoothing-]{smoothing}
\externaldocument[modules-]{modules}
\externaldocument[sites-modules-]{sites-modules}
\externaldocument[injectives-]{injectives}
\externaldocument[cohomology-]{cohomology}
\externaldocument[sites-cohomology-]{sites-cohomology}
\externaldocument[dga-]{dga}
\externaldocument[dpa-]{dpa}
\externaldocument[sdga-]{sdga}
\externaldocument[hypercovering-]{hypercovering}
\externaldocument[schemes-]{schemes}
\externaldocument[constructions-]{constructions}
\externaldocument[properties-]{properties}
\externaldocument[morphisms-]{morphisms}
\externaldocument[coherent-]{coherent}
\externaldocument[divisors-]{divisors}
\externaldocument[limits-]{limits}
\externaldocument[varieties-]{varieties}
\externaldocument[topologies-]{topologies}
\externaldocument[descent-]{descent}
\externaldocument[perfect-]{perfect}
\externaldocument[more-morphisms-]{more-morphisms}
\externaldocument[flat-]{flat}
\externaldocument[groupoids-]{groupoids}
\externaldocument[more-groupoids-]{more-groupoids}
\externaldocument[etale-]{etale}
\externaldocument[chow-]{chow}
\externaldocument[intersection-]{intersection}
\externaldocument[pic-]{pic}
\externaldocument[weil-]{weil}
\externaldocument[adequate-]{adequate}
\externaldocument[dualizing-]{dualizing}
\externaldocument[duality-]{duality}
\externaldocument[discriminant-]{discriminant}
\externaldocument[derham-]{derham}
\externaldocument[local-cohomology-]{local-cohomology}
\externaldocument[algebraization-]{algebraization}
\externaldocument[curves-]{curves}
\externaldocument[resolve-]{resolve}
\externaldocument[models-]{models}
\externaldocument[functors-]{functors}
\externaldocument[equiv-]{equiv}
\externaldocument[pione-]{pione}
\externaldocument[etale-cohomology-]{etale-cohomology}
\externaldocument[proetale-]{proetale}
\externaldocument[relative-cycles-]{relative-cycles}
\externaldocument[more-etale-]{more-etale}
\externaldocument[trace-]{trace}
\externaldocument[crystalline-]{crystalline}
\externaldocument[spaces-]{spaces}
\externaldocument[spaces-properties-]{spaces-properties}
\externaldocument[spaces-morphisms-]{spaces-morphisms}
\externaldocument[decent-spaces-]{decent-spaces}
\externaldocument[spaces-cohomology-]{spaces-cohomology}
\externaldocument[spaces-limits-]{spaces-limits}
\externaldocument[spaces-divisors-]{spaces-divisors}
\externaldocument[spaces-over-fields-]{spaces-over-fields}
\externaldocument[spaces-topologies-]{spaces-topologies}
\externaldocument[spaces-descent-]{spaces-descent}
\externaldocument[spaces-perfect-]{spaces-perfect}
\externaldocument[spaces-more-morphisms-]{spaces-more-morphisms}
\externaldocument[spaces-flat-]{spaces-flat}
\externaldocument[spaces-groupoids-]{spaces-groupoids}
\externaldocument[spaces-more-groupoids-]{spaces-more-groupoids}
\externaldocument[bootstrap-]{bootstrap}
\externaldocument[spaces-pushouts-]{spaces-pushouts}
\externaldocument[spaces-chow-]{spaces-chow}
\externaldocument[groupoids-quotients-]{groupoids-quotients}
\externaldocument[spaces-more-cohomology-]{spaces-more-cohomology}
\externaldocument[spaces-simplicial-]{spaces-simplicial}
\externaldocument[spaces-duality-]{spaces-duality}
\externaldocument[formal-spaces-]{formal-spaces}
\externaldocument[restricted-]{restricted}
\externaldocument[spaces-resolve-]{spaces-resolve}
\externaldocument[formal-defos-]{formal-defos}
\externaldocument[defos-]{defos}
\externaldocument[cotangent-]{cotangent}
\externaldocument[examples-defos-]{examples-defos}
\externaldocument[algebraic-]{algebraic}
\externaldocument[examples-stacks-]{examples-stacks}
\externaldocument[stacks-sheaves-]{stacks-sheaves}
\externaldocument[criteria-]{criteria}
\externaldocument[artin-]{artin}
\externaldocument[quot-]{quot}
\externaldocument[stacks-properties-]{stacks-properties}
\externaldocument[stacks-morphisms-]{stacks-morphisms}
\externaldocument[stacks-limits-]{stacks-limits}
\externaldocument[stacks-cohomology-]{stacks-cohomology}
\externaldocument[stacks-perfect-]{stacks-perfect}
\externaldocument[stacks-introduction-]{stacks-introduction}
\externaldocument[stacks-more-morphisms-]{stacks-more-morphisms}
\externaldocument[stacks-geometry-]{stacks-geometry}
\externaldocument[moduli-]{moduli}
\externaldocument[moduli-curves-]{moduli-curves}
\externaldocument[examples-]{examples}
\externaldocument[exercises-]{exercises}
\externaldocument[guide-]{guide}
\externaldocument[desirables-]{desirables}
\externaldocument[coding-]{coding}
\externaldocument[obsolete-]{obsolete}
\externaldocument[fdl-]{fdl}
\externaldocument[index-]{index}

% Heap Project documents
\externaldocument[k3-]{k3}
\externaldocument[torelli-]{torelli}
\externaldocument[linear-algebra-]{linear-algebra}
\externaldocument[groups-]{groups}
\externaldocument[complex-analysis-]{complex-analysis}
\externaldocument[functional-analysis-]{functional-analysis}
\externaldocument[categories-]{categories}
\externaldocument[commutative-algebra-]{commutative-algebra}
\externaldocument[ringed-spaces-]{ringed-spaces}
\externaldocument[homological-algebra-]{homological-algebra}
\externaldocument[lie-algebras-]{lie-algebras}
\externaldocument[representation-theory-]{representation-theory}
\externaldocument[smooth-manifolds-]{smooth-manifolds}
\externaldocument[differential-topology-]{differential-topology}
\externaldocument[algebraic-topology-]{algebraic-topology}
\externaldocument[differential-geometry-]{differential-geometry}
\externaldocument[algebraic-geometry-]{algebraic-geometry}
\externaldocument[symplectic-geometry-]{symplectic-geometry}
\externaldocument[complex-geometry-]{complex-geometry}
\externaldocument[hodge-theory-]{hodge-theory}
\externaldocument[%
variations-of-hodge-structures-]{%
variations-of-hodge-structures}
\externaldocument[%
homological-mirror-symmetry-]{%
homological-mirror-symmetry}
\maybeexternaldocument{deformations-}{deformations}
\externaldocument[vertex-algebras-]{vertex-algebras}
\externaldocument[springer-]{springer}
\externaldocument[infty-categories-]{infty-categories}
\externaldocument[seminars-differential-geometry-]{%
seminars-differential-geometry}
\externaldocument[seminars-complex-geometry-]{%
seminars-complex-geometry}
\externaldocument[seminars-algebraic-geometry-]{%
seminars-algebraic-geometry}
\externaldocument[seminars-analysis-pde-]{%
seminars-analysis-pde}
\externaldocument[surfaces-]{surfaces}
\externaldocument[birational-maps-commutative-algebra-]{birational-maps-commutative-algebra}
\externaldocument[cremona-transformations-]{cremona-transformations}
\externaldocument[derived-categories-of-sheaves-]{derived-categories-of-sheaves}
\externaldocument[geometric-invariant-theory-]{geometric-invariant-theory}
\externaldocument[geometric-stability-conditions-and-group-actions-]{geometric-stability-conditions-and-group-actions}
\externaldocument[moduli-spaces-of-sheaves-]{moduli-spaces-of-sheaves}
\externaldocument[bridgeland-stability-general-theory-]{bridgeland-stability-general-theory}
\externaldocument[hilbert-schemes-]{hilbert-schemes}
\externaldocument[noncommutative-abelian-surfaces-and-kummer-type-hyperkahler-varieties-]{noncommutative-abelian-surfaces-and-kummer-type-hyperkahler-varieties}
\externaldocument[mumford-tate-groups-in-hodge-theory-]{mumford-tate-groups-in-hodge-theory}
\externaldocument[hodge-birational-atoms-2026-]{hodge-birational-atoms-2026}
\externaldocument[xxii-egd-2026-]{%
xxii-egd-2026}
\externaldocument[pde-]{pde}
\externaldocument[probability-]{probability}

\externaldocument[linguistics-]{linguistics}

\externaldocument[%
16th-alga-meeting-2026-]{%
16th-alga-meeting-2026}

\externaldocument[%
iv-jornadas-lefschetz-2026-]{%
iv-jornadas-lefschetz-2026}

\externaldocument[reading-the-masters-]{%
reading-the-masters}

\externaldocument[real-analysis-]{%
real-analysis}

\externaldocument[optimal-transport-]{%
optimal-transport}

\externaldocument[history-]{history}

\externaldocument[seiberg-witten-]{seiberg-witten}

\externaldocument[history-philosophy-biology-]{history-philosophy-biology}

\externaldocument[classical-mechanics-]{classical-mechanics}
\externaldocument[electromagnetism-]{electromagnetism}
\externaldocument[waves-and-optics-]{waves-and-optics}
\externaldocument[wif-]{waves-in-fluids-iv}
\externaldocument[thermal-physics-]{thermal-physics}
\externaldocument[quantum-mechanics-]{quantum-mechanics}
\externaldocument[special-relativity-]{special-relativity}
\externaldocument[physics-]{physics}

\externaldocument[chemistry-]{chemistry}

\externaldocument[molecular-cell-biology-]{molecular-cell-biology}

\externaldocument[microbiology-]{microbiology}

\externaldocument[protistology-]{protistology}

\externaldocument[mycology-]{mycology}

\externaldocument[botany-]{botany}

\externaldocument[zoology-]{zoology}

\externaldocument[animal-histology-]{animal-histology}

\externaldocument[animal-physiology-]{animal-physiology}

\externaldocument[systematics-biogeography-]{systematics-biogeography}

\externaldocument[ecology-natural-resources-]{ecology-natural-resources}

\externaldocument[earth-sciences-]{earth-sciences}

\externaldocument[biostatistics-]{biostatistics}

\externaldocument[mathematical-systems-biology-]{mathematical-systems-biology}

\externaldocument[cimpa-el-salvador-2024-]{cimpa-el-salvador-2024}

\externaldocument[anthropology-]{anthropology}

\externaldocument[yang-mills-theory-]{yang-mills-theory}

\externaldocument[seminars-algebra-]{seminars-algebra}

\externaldocument[seminars-physics-]{seminars-physics}

\externaldocument[seminars-computing-]{seminars-computing}

\externaldocument[oktobermat-2026-]{oktobermat-2026}

\externaldocument[saet-2026-]{saet-2026}

\externaldocument[seminars-symplectic-geometry-]{seminars-symplectic-geometry}

\externaldocument[seminars-dynamics-]{seminars-dynamics}

\externaldocument[oktobermat-2024-]{oktobermat-2024}

\externaldocument[quantum-groups-]{quantum-groups}

\externaldocument[legal-2027-]{legal-2027}

\externaldocument[derived-categories-]{derived-categories}

% Theorem environments.
%
\theoremstyle{plain}
\newtheorem{theorem}[subsection]{Theorem}
\newtheorem{proposition}[subsection]{Proposition}
\newtheorem{lemma}[subsection]{Lemma}
\newtheorem{corollary}[subsection]{Corollary}

\theoremstyle{definition}
\newtheorem{definition}[subsection]{Definition}
\newtheorem{example}[subsection]{Example}
\newtheorem{exercise}[subsection]{Exercise}
\newtheorem{situation}[subsection]{Situation}

\theoremstyle{remark}
\newtheorem{remark}[subsection]{Remark}
\newtheorem{remarks}[subsection]{Remarks}

\numberwithin{equation}{subsection}

% Macros
%
\def\lim{\mathop{\mathrm{lim}}\nolimits}
\def\colim{\mathop{\mathrm{colim}}\nolimits}
\def\Spec{\mathop{\mathrm{Spec}}}
\def\Hom{\mathop{\mathrm{Hom}}\nolimits}
\def\Ext{\mathop{\mathrm{Ext}}\nolimits}
\def\SheafHom{\mathop{\mathcal{H}\!\mathit{om}}\nolimits}
\def\SheafExt{\mathop{\mathcal{E}\!\mathit{xt}}\nolimits}
\def\Sch{\mathit{Sch}}
\def\Mor{\mathop{\mathrm{Mor}}\nolimits}
\def\Ob{\mathop{\mathrm{Ob}}\nolimits}
\def\Sh{\mathop{\mathit{Sh}}\nolimits}
\def\NL{\mathop{N\!L}\nolimits}
\def\CH{\mathop{\mathrm{CH}}\nolimits}
\def\proetale{{pro\text{-}\acute{e}tale}}
\def\etale{{\acute{e}tale}}
\def\QCoh{\mathit{QCoh}}
\def\Ker{\mathop{\mathrm{Ker}}}
\def\Im{\mathop{\mathrm{Im}}}
\def\Coker{\mathop{\mathrm{Coker}}}
\def\Coim{\mathop{\mathrm{Coim}}}

% Boxtimes
%
\DeclareMathSymbol{\boxtimes}{\mathbin}{AMSa}{"02}

%
% Macros for moduli stacks/spaces
%
\def\QCohstack{\mathcal{QC}\!\mathit{oh}}
\def\Cohstack{\mathcal{C}\!\mathit{oh}}
\def\Spacesstack{\mathcal{S}\!\mathit{paces}}
\def\Quotfunctor{\mathrm{Quot}}
\def\Hilbfunctor{\mathrm{Hilb}}
\def\Curvesstack{\mathcal{C}\!\mathit{urves}}
\def\Polarizedstack{\mathcal{P}\!\mathit{olarized}}
\def\Complexesstack{\mathcal{C}\!\mathit{omplexes}}
% \Pic is the operator that assigns to X its picard group, usage \Pic(X)
% \Picardstack_{X/B} denotes the Picard stack of X over B
% \Picardfunctor_{X/B} denotes the Picard functor of X over B
\def\Pic{\mathop{\mathrm{Pic}}\nolimits}
\def\Picardstack{\mathcal{P}\!\mathit{ic}}
\def\Picardfunctor{\mathrm{Pic}}
\def\Deformationcategory{\mathcal{D}\!\mathit{ef}}


\begin{document}

\title{Seminars in Symplectic Geometry}
\maketitle

\phantomsection
\label{section-phantom}

\tableofcontents

\section{Magnetic monopoles according to
Hitchin}
\label{section-magnetic-monopoles-hitchin}
%% speaker: Graham Smith, PUC-Rio
%% place: Seminário das Sextas, PUC-Rio
%% date: February 8, 2024

\noindent
Graham Smith, PUC-Rio.
Seminário das Sextas, PUC-Rio.
8 of February 2024.

\begin{definition}
\label{definition-mono-magnetic-monopole}
A \emph{magnetic monopole} consists of 3
elements: $(E,\nabla,\Phi)$, a bundle, a
connection, and the Higgs field.
\end{definition}

And the Bogomolny equations are:

\begin{lemma}[Bogomolny equations]
\label{lemma-mono-bogomolny}
$\nabla$ is anti self dual iff $*
F^{\nabla^1}=\nabla^1\Phi_0$.
\end{lemma}

\medskip\noindent
{\bf What's going on} (I arrived in late).
ChatGPT says $F^\nabla$ is the curvature of
this connection on this principal bundle. So
Bogomolny equations are equations on the
curvature. Like what? {\bf Like Maxwell's
equations!}

Here's some more ChatGPT on what's about to
come in this talk:

\medskip\noindent
{\bf Magnetic Monopoles and Twistors.} A
magnetic monopole consists of a vector bundle
$E$, a connection $\nabla$, and a Higgs field
$\Phi$. The Bogomolny equation, $*F^\nabla =
\nabla \Phi$, describes monopole solutions
balancing curvature and the Higgs field. The
Higgs field $\Phi$ can be interpreted as a
difference between two connections or as a
section of $\operatorname{End}(E)$.

Twistor theory provides a complex-geometric
approach to gauge fields. The twistor space
of $\mathbb{R}^3$ can be described in terms
of oriented lines, isotropic
$\mathbb{C}$-lines, or null
$\mathbb{C}$-planes in $\mathbb{C}^3$. For a
$\mathsf{U}(2)$-connection, one constructs a
holomorphic vector bundle $\tilde{E} \to X$
over the complex manifold $X = TS^2$,
encoding monopoles in terms of holomorphic
structures.

\begin{remark}[Misha]
\label{remark-mono-misha}
$\Phi$ is sections of endomorphisms of the
group which is $G$.
\end{remark}

\medskip\noindent
{\bf Question.} What is the Higgs field? It
is a \emph{difference} between two
connections.

\medskip\noindent
{\bf Question.} What are twistors?

\medskip\noindent
{\bf Twistors in $\mathbb{R}^3$.}
Embed $\mathbb{R}^3 \hookrightarrow
\mathbb{C}^3$. Extender the inner product of
$\mathbb{R}^3$ $\left\langle{} \cdot,\cdot
\right\rangle{}$ \emph{blinearly} to
$\mathbb{C}^3$. (It's not hermitian, it's a
bilinear form.)

Twistor space of $\mathbb{R}^3$:

\begin{enumerate}
\item Oriented lines in $\mathbb{R}^3$.
\item Null $\mathbb{C}$-lines in
$\mathbb{C}^3$ (isotropic).
\item Null $\mathbb{C}$-planes in
$\mathbb{C}^3$.
\end{enumerate}
The three are equivalent (we have shown that
in the seminar).

It turns out that this space can be
parametrized by $TS^3$ as follows: the
oriented line $L^+$ goes to the pair $(x,y)$,
where $x$ is the unit tangent and $y$ is the
closest point.
$$
x \in S^2 \qquad y \perp x \iff y \in T_x
S^2.
$$
Now we have a complex integrable structre
becuase $S^2=\hat{\mathbb{C}}$.

Now

suppose $(E^2, \nabla, \Phi)$. Fix the group
to be $\mathsf{U}(2)$ (it's a
$\mathsf{U}(2)$-connection). Define a bundle
$\tilde{E}$ over $X$ by
$$
\begin{aligned}
\tilde{E}_x &= \{ \sigma: L_x \to R:
\nabla_T\sigma - u\Phi \sigma=0\} \\
&= \{\sigma:L_x \to E: \tilde{\nabla}_{(T+i
e_0)}\sigma=0\}
\end{aligned}
$$
where $\tilde{\nabla}$ is the Yang-Mills
connection.

Now take the complex manifold $X= TS^2$, so
$(E^2,\nabla,\Phi)$ and produce the bundle
$\tilde{ E} \to X$. Then you may actually put
a holomorphic structure on $\tilde{E}$,
making into a holomorphic vector bundle over
$X$.

Now we write: $X$ is the set of oriented
lines in $\mathbb{R}^3$, and
$$
\begin{aligned}
T_{(x,y)}TS^2 &= \{(u,v): \left\langle{} u,v
\right\rangle{}=0, \left\langle{} u,y
\right\rangle{}+\left\langle{} x,v
\right\rangle{}=0\} \\
V_x &= \{(0,v) : \left\langle{} x,v
\right\rangle{}=0\} \\
J(0,v) &= (0,x \times v) \\
H_{(x,y)}TS^2 &= \{(u,-\left\langle{} u,y
\right\rangle{}x:\left\langle{} u,x
\right\rangle{}=0\} \\
\hat{u} &:=(u,-\left\langle{} u,y
\right\rangle{}x) \\
J\hat{u} &= \widehat{x \times u}
\end{aligned}
$$

And then: Jac fields orthogonal to $L^+$ is 4
dimensional.
$$
J_{\xi}=T \times \xi
$$
Constant curvature $\to$ integrable.

$$
\begin{aligned}
\gamma_{x,y}(t) &= y+t x \\
(u,v- &\left\langle{} u,y \right\rangle{}x \\
\xi(t) &= v-\left\langle{} u,y
\right\rangle{}x+tu
\end{aligned}
$$
$G$ lie group, $H\subseteq G$
$$
\begin{aligned}
\mathfrak{g} &= \mathfrak{h} \oplus
\mathfrak{p} \\
[\mathfrak{g},\mathfrak{p}] &\subseteq
\mathfrak{p} \qquad
[\mathfrak{p},\mathfrak{p}] \subseteq
\mathfrak{g}
\end{aligned}
$$
(Polarization.) $\mathfrak{p}$ id. with
$T(G/H)$ fwy
$\operatorname{Ad}_\mathfrak{g}$-invariant
form on $\mathfrak{p}$ generates a parallel
form on $G/H$.

$$
\begin{aligned}
G &= \mathsf{SO}(3)\ltimes \mathbb{R}^3 \\
H &= \mathsf{SO}(2) \times \mathbb{R} \\
\mathfrak{p} &= 4, J, \nabla J=0
\end{aligned}
$$

Here's ChatGPT:

\medskip\noindent
{\bf Twistor Space and Structure.} The
twistor space of $\mathbb{R}^3$ is given by
$X = TS^2$, which parametrizes oriented lines
in $\mathbb{R}^3$. The tangent space at
$(x,y) \in TS^2$ decomposes as
$$
T_{(x,y)}TS^2 = \{(u,v) : \left\langle{} u,v
\right\rangle{} = 0, \left\langle{} u,y
\right\rangle{} + \left\langle{} x,v
\right\rangle{} = 0 \}.
$$
The vertical and horizontal subspaces are
defined as
$$
V_x = \{(0,v) : \left\langle{} x,v
\right\rangle{} = 0\}, \quad H_{(x,y)}TS^2 =
\{(u,-\left\langle{} u,y \right\rangle{} x) :
\left\langle{} u,x \right\rangle{} = 0\}.
$$
An almost complex structure $J$ is introduced
by
$$
J(0,v) = (0, x \times v), \quad J\hat{u} =
\widehat{x \times u}.
$$

\medskip\noindent
{\bf Jacobi Fields and Integrability.} The
Jacobi fields orthogonal to an oriented line
$L^+$ form a 4-dimensional space, where
$$
J_{\xi} = T \times \xi.
$$
Constant curvature ensures integrability.

The geodesic equation for a curve
$\gamma_{x,y}(t) = y + tx$ leads to
$$
\xi(t) = v - \left\langle{} u,y
\right\rangle{} x + tu.
$$

\medskip\noindent
{\bf Lie Groups and Homogeneous Structure.}
Consider a Lie group $G$ with a subgroup $H
\subseteq G$, and the Lie algebra
decomposition
$$
\mathfrak{g} = \mathfrak{h} \oplus
\mathfrak{p}, \quad
[\mathfrak{g},\mathfrak{p}] \subseteq
\mathfrak{p}, \quad
[\mathfrak{p},\mathfrak{p}] \subseteq
\mathfrak{g}.
$$
Since $\mathfrak{p}$ is identified with
$T(G/H)$, an
$\operatorname{Ad}_{\mathfrak{g}}$-invariant
form on $\mathfrak{p}$ induces a parallel
form on $G/H$.

For the twistor space,
$$
G = \mathsf{SO}(3) \ltimes \mathbb{R}^3,
\quad H = \mathsf{SO}(2) \times \mathbb{R}.
$$
The space $\mathfrak{p}$ has dimension 4, and
the almost complex structure $J$ satisfies
$\nabla J = 0$, ensuring integrability.

\medskip\noindent
{\bf Scattering transform.}
Suppose you have a Lie group $G$, and two Lie
subgroups $H_1, H_2 \subseteq G$. Now take
$$
X:=G/H_1, \qquad X_2=G/H_2
$$
[Figure in the original notes: commutative
diagram with $G$ on top, mapping to $X_1$ by
$\pi_1$ and to $X_2$ by $\pi_2$.]

{\bf Radon transform}:
$$
R: C^\infty(X_1) \longrightarrow
C^\infty(X_2)
$$
So
$$
\begin{aligned}
R &= \pi_{2 *}\pi_1^* \\
G &= \mathsf{SO}(3)\ltimes \mathbb{R}^3 \\
H_1 &= \mathsf{SO}(3) \\
H_2 &= S^1 \times \mathbb{R}
\end{aligned}
$$
Here $G$ is the group of isometries of
$\mathbb{R}^3$, $H_1$ is the stabilizer of a
point (``Stab point'' in the original notes)
and $H_2$ is the stabilizer of a line
(``Stable line''). Now we can put a little
zero on the continuous functions (that means
they tend to zero?):
$$
R: C_0^\infty(X_1) \longrightarrow
C^\infty(X_2)
$$
$$
\begin{aligned}
R &= \pi_{2*}\pi_1^* \\
\pi_1^* (f) &= f \circ \pi_1 \\
\pi_{2*}f &= \int_HfdV\frac{1}{2\pi} \\
\hat{f}(L) &= \int_Lfde \\
F &= \pi_{2*}\mu_{\phi}\pi_1^*
\end{aligned}
$$

[Figure in the original notes: commutative
diagram with $G=\mathsf{SO}(3)\ltimes
\mathbb{R}^3$ on top, mapping to
$\mathbb{R}^3$ by $\pi_1$ and to $X$ by
$\pi_2$, and below $G$ the triple
$(\pi_1^*E,\pi_1^*\nabla,\pi^*_1\phi)$ with a
vertical arrow up to $G$.]

So a global section
$$
(\pi_1^*E_1|_{F},
\pi_1^*\nabla|_{F},
\pi_1^*\phi|_{F}
$$

\medskip\noindent
{\bf Summary of the Scattering Transform.}
Given a Lie group $G$ and two subgroups $H_1,
H_2 \subseteq G$, we define the homogeneous
spaces
$$
X_1 = G/H_1, \quad X_2 = G/H_2.
$$
The \emph{Radon transform} is a mapping
$$
R: C^\infty_0(X_1) \to C^\infty(X_2),
$$
which factors through the pullback $\pi_1^*$
and pushforward $\pi_{2*}$:
$$
R = \pi_{2*} \pi_1^*.
$$

For $G = \mathsf{SO}(3) \ltimes \mathbb{R}^3$
(the isometry group of $\mathbb{R}^3$), the
subgroups correspond to:
$$
H_1 = \mathsf{SO}(3), \quad H_2 = S^1 \times
\mathbb{R}.
$$
Here $H_1$ stabilizes a point and $H_2$
stabilizes a line. The transform integrates
functions along fibers:
$$
\pi_1^* f = f \circ \pi_1, \quad \pi_{2*} f =
\int_H f dV.
$$
A section of a lifted bundle $(\pi_1^*E,
\pi_1^*\nabla, \pi_1^*\phi)$ over $G$
descends through $\pi_2$, encoding the
scattering transform in terms of differential
geometric data.

\medskip\noindent
{\bf Editorial note.} In porting, some
labels that stood inside formulas in the
original notes were moved to prose: the
word ``Polarization'', the descriptions of
$G$, $H_1$, $H_2$ (isometries, stabilizer of
a point, of a line) and the sentence that
$X$ is the set of oriented lines. The two
commutative diagrams are replaced by
bracketed descriptions. Unbalanced
parentheses and the unfinished formula
for the global section are as in the
original notes, not repaired.

\section{Geometric quantization on K\"ahler
manifolds, Berezin-Toeplitz, coherent
states}
\label{section-geometric-quantization-kahler}
%% speaker: Bruno Suassuna, PUC-Rio
%% place: Seminário das Sextas, PUC-Rio
%% date: October 18, 2024

\noindent
Bruno Suassuna, PUC-Rio.
Seminário das Sextas, PUC-Rio.
18 October 2024.

\medskip
The idea of geometric quantization is to
start with a symplectic manifold $(M,\omega)$
and for some geometric structure that you
have on $M$, like a line bundle with a
hermitian metric, which is the geometric
quantization data.

So on every point $x \in M$ you have a copy
of $\mathbb{C}$, and you put a hermitian
metric $h_x:\mathcal{L}_x\otimes
\bar{\mathcal{L}}_x\to \mathbb{C}$.

Now the second part of the geometric
quantization data is a connection $\nabla$ on
$\mathcal{L}$ that preserves $h$. Recall that
this means that
$$
\nabla :\mathcal{L}\to \mathcal{L}\otimes
\Omega^{1}(M)
$$
such that
\begin{itemize}
\item $\forall v\in\mathcal{C}^\infty(TM)$,
$\nabla_v:\mathcal{L}\to \mathcal{L}$ is
$\mathbb{C}$-linear
\item Leibniz in the sense that
$\nabla_v(fs)=df(v)s+f \nabla_vs$.
\end{itemize}

\medskip\noindent
{\bf Property} (Bohr-Sommerfield).
$\frac{2\pi}{i}R_\nabla =\omega$, where
$(E,\nabla )$ is a vector bundle with
connection we define $R_\nabla
\in\Omega^{2}_M \otimes
\operatorname{End}(E)$ (the curvature).

\begin{definition}
\label{definition-suassuna-quantizable}
$M$ is \emph{quantizable} when it admits the
data above satisfyint the property
Bohr-Sommerfield.

We call $L^2(M,\mathcal{L})$ a
\emph{pre-quantum Hilbert space}.
\end{definition}

\medskip\noindent
{\bf Question.}
What is it that makes it better in Kähler
manifolds? (Dani: I thought in general the
function space didn't have to be the sections
of a bundle)

\begin{remark}[Sergey]
\label{remark-suassuna-too-big}
When $\mathbb{R}^{2n}$ is the phase space and
$\mathcal{L}$ is the trivial bundle then
$L^2( \mathbb{R}^{2} )$ is \emph{too big!}
\end{remark}

\begin{definition}
\label{definition-suassuna-kahler}
A \emph{K\"ahler manifold} is $(M,
\omega,g,I)$ where $\omega$ is a symplectic
structure, $g$ a riemannian structure and $I$
a complex structure, and they are all
compatible and $I$ is integrable. So remember
that compatibility is for example that
$g(u,v)=\omega(Iu,v)$ and $g$ is
$I$-invariant. On the other hand,
integrability is a non-trivial PDE and gives
holomorphic local coordinates.
\end{definition}

{\bf Starting again,}\hspace{.5em} take a
K\"ahler manifold $M$ and define a
pre-quantum line bundle so $\mathcal{L}$
holomorphic with the Chern connection
$\nabla$. Also suppose that $\nabla$
satisfies the Bohr-Sommerfield property.

\begin{remark}[Kodaira theorem]
\label{remark-suassuna-kodaira}
The K\"ahler class (cohomology class of
symplectic form) is the first Chern class,
i.e., $[\omega]=c_1(\ell)\in
H^{2}(M,\mathbb{Z})$ iff $\mathcal{L}$ is
ample.
\end{remark}

\begin{definition}
\label{definition-suassuna-hilbert-space}
$$
\mathcal{H}_m=H^{0}(X,\mathcal{L}^{\otimes
m})
$$
so the sections of that bundle.
\end{definition}

You see hare making the space of sections
smaller:
$$
\mathcal{H}_m=H^{0}(X,\mathcal{L}^{\otimes
m})\hookrightarrow L^2(M,\mathcal{L}^{\otimes
m})
$$
And if $M$ is not compact, take
$$
H^{0}(M,\mathcal{L}^{\otimes m})\cap L_2
$$

\medskip\noindent
{\bf Question} (Altan).
What about the measure here?

\medskip\noindent
{\bf Answer} (Bruno).
Related to Liouville measure; a clever choice
of $h$. So maybe
$$
\left\langle{} s_1,s_2 \right\rangle{}
=\int_{M}h\left( s_1(x),s_2(x) \right)
\frac{\omega^d}{d!}
$$
And now $d$ is what in Altan's talk was
$n$\ldots{} or was it $m$?

\begin{remark}
\label{remark-suassuna-altan}
So in the case of Altan's talk
$$
\mathcal{H}=\left\{
f:\mathbb{C}^{n}\overset{\operatorname{entire}}{\longrightarrow}
\mathbb{C}:\int_{\mathbb{C}^{n}}|f|^2
\operatorname{exp}(-|z|^2)
d \lambda(z) \right\}
$$
so notice that the $\operatorname{exp}$ term
is stopping the space to be trivial because
bounded holomorphic functions are constant
right?
\end{remark}

$\mathsf{OK}$ back to the quantization. Given
$f \in\mathcal{C}^\infty(M)$ define $A_f
\in\operatorname{End}(\mathcal{H}_m)$\ldots{}

\medskip\noindent
{\bf Berezin-Toeplitz.}
$$
T_f(S)=\Pi_m(fS)
$$
$\Pi_m$ orthogonal projection to
$H^{0}(M,\mathcal{L}^{\otimes m})$.

\medskip\noindent
{\bf Kostant-Sorian pre-quantum operators.}
$f$ quantizable means
$Q_f(s):=\nabla_{X_f}s-2\pi i fs$. Here $X_f$
is the hamiltonian vector field of $f$ with
respect to $\omega$.

So that's an operator associated to a
function---quantization!

And we \emph{wish} that if $s$ is holomorphic
then so is $Q_f(s)$.

\begin{remark}[Sergey]
\label{remark-suassuna-wish-list}
Also in the wish list is some property of
brackets right?
\end{remark}

Right so we are doing
$$
\begin{aligned}
\mathcal{C}^\infty(M) &\longrightarrow
\operatorname{End}(V) \\
f &\longmapsto \hat{f}
\end{aligned}
$$
and we want
$$
[\hat{f},\hat{g}]=i\hbar \widehat{\{f,g\}}
\qquad \hat{1}=\operatorname{Id}, \qquad
\widehat{\varphi\circ f}=\varphi(\hat{f})
$$

\medskip\noindent
{\bf Question} (Dani).
So why is it impossible that the wishes
became true? Is that a difficult theorem?

Now
$$
L_f(s)=\Pi_m(Q_f(s))
$$
and

\medskip\noindent
{\bf Tuyman's lemma.}
$iQ_f=T_f-\perp_{2m}\Delta f$

\medskip\noindent
{\bf Coherent state quantization.}
$p \in \mathcal{L}_x^{\otimes m}\setminus
\{0\} \rightsquigarrow e_p \in
\mathcal{H}_m$. So
$$
s \in H^{0}(X,\mathcal{L}^{\otimes
m})\longmapsto
s(x)\in\mathcal{L}_x=\lambda_p(s)\cdot p
$$
Now a \emph{coherent state} is defined via
Riesz' representation theorem so $\lambda_p
\in\mathcal{H}_m^*$ such that
$$
\lambda_p(s)=\left\langle{} e_p,s
\right\rangle{}
$$

\medskip\noindent
{\bf Question} (Dani).
Why are coherent states so important?

They give us a rational map which is
essentially Kodaira map (the map defined by
any line bundle)
$$
\operatorname{coh}:M\dashrightarrow
\mathbb{P}(\overline{\mathcal{H}}_m)
$$
So given a bounded operator $A\in
B(\mathcal{H}_m)$ we define
$$
\hat{A}(x)=\frac{\left\langle{} Ae_p,e_p
\right\rangle{}}{\|e_p\|^2} \qquad x\in M
$$
called the \emph{covariant symbol}. (That is
not the same hat than the hat of the
quantization!) We have a map $B
(\mathcal{H}_m)\to \mathcal{C}^\infty(M)$.

Now we try to invert this map: which
functions are the covariant symbols of some
bounded operators?

\begin{remark}[Sergey]
\label{remark-suassuna-kernel}
This is like in altan talk, look for the
kernel!
\end{remark}

Here's two formulas to go in the other
direction:

$$
(As)(x)=\int_{M}h_y(s(y),s(y))
\hat{A}(x,y)\frac{\omega^m}{n!}
$$
$$
\operatorname{Tr}(A)=
\int_{M}\hat{A}(x)\theta(x)\frac{\omega^n}{n!}
$$
and that's a nice formula wich has to do with
Rawsley, $\theta(x)=|q|^2\cdot\|e_q\|^2$.

\begin{remark}[Dani]
\label{remark-suassuna-equivalence}
So it looks like we want an equivalence of
operators and functions.
\end{remark}

\medskip\noindent
{\bf Bondermann-Meinken-Schlihenmeier.}
$$
\|f\|_\infty-\frac{c}{m}\leq \|T_f^{m}\|\leq
\|f\|_\infty
$$
$$
\|m[T_f^{(m)}, T_y^{(m)}
]iT_{\{f,g\}}^{(m)}\|= \mathcal{O}(m^{-1})
$$
as $n\to \infty$.

\medskip\noindent
{\bf Editorial note.} Repairs made in
porting: a missing bracket in
$[\hat{f},\hat{g}]$ and a missing closing
parenthesis in
$\mathbb{P}(\overline{\mathcal{H}}_m)$
were added. The last formula (commutator
estimate) is as in the original notes and
looks garbled (the names, the symbol $y$
and the operation between the commutator
and $iT_{\{f,g\}}$); it is not repaired.

\section{Symplectic and contact nature of
Riemannian geometry}
\label{section-symplectic-contact-riemannian}
%% speaker: Graham Smith, PUC-Rio
%% place: Seminário das Sextas, PUC-Rio
%% date: November 1 and 8, 2024

\noindent
Graham Smith, PUC-Rio.
Seminário das Sextas, PUC-Rio.
1st and 8th of November 2024.

\medskip
Surfaces of constant extrinsic curvature $k$.

\begin{center}
\begin{tabular}{c |c| c| c}
& $\mathbb{H}^2$ & $\mathbb{R}^3$ & $S^3$ \\
\hline
$k>1$ & $S^3$ & $S^2$ & $S^2$ \\
$k=1$ & Horospheres, horocycles & & \\
$0<k<1$ & ess. param.,
$\mathcal{H}(\mathbb{D})\sqcup
\mathcal{H}(\mathbb{C})\setminus
\{\mathbb{C}\}$
\end{tabular}
\end{center}

\begin{remark}
\label{remark-smith-curvatures}
Intrinsic curvature = extrinsic curvature +
sectional curvature of the ambient space.
\end{remark}

$X$ a 3-manifold, take $x \in X$ and $v_x \in
TX$. The Levi-Civita connection provides
$$
\begin{aligned}
T_vTX &\cong H_vTX\oplus V_xTX \\
&\cong T_xX \oplus T_xX
\end{aligned}
$$
where the first copy of $T_xX$ is the
horizontal part and the second is the
vertical part.
\begin{itemize}
\item Saski metric
$$
\left\langle{} (\xi,\mu),(\xi',\mu')
\right\rangle{}=\left\langle{} \xi,\xi'
\right\rangle{}+\left\langle{} \mu,\mu'
\right\rangle{}
$$
\item Symplectic form
$$
\omega\Big((\xi,\mu),(\xi',\mu')\Big)=
\left\langle{}
\xi,\mu' \right\rangle{}-\left\langle{}
\xi',\mu \right\rangle{}
$$
which we may pull back to $X$ via the musical
isomorphism to obtain the \emph{Saski
symplectic form} of $T^*X$
$\omega=\flat^*\omega_{\operatorname{st}}$.
\item $$m\Big((\xi,\mu),(\xi',\mu')\Big)$$
\item There is a complex structure $I$
\item And a quadratic form
$$
m=\begin{pmatrix}0&I_0\\ I_0&0\end{pmatrix}
$$
\end{itemize}
And a contact bundle pulling the luiville
form from $T^*X$. And then there is a another
complex structure $J$. So define $K:=I J$,
which gives a {\bf hyperkähler contact
structure} on the unit tangent bundle.

\medskip\noindent
{\bf Upshot.}
We have created a pseudoholomorphic curve
from a $k$-surface by taking the unit tangent
normal. That is, consider the normal map as
an embedding of our surface in the unit
tangent bundle.

\begin{theorem}[Jurgens]
\label{theorem-smith-jurgens}
$\Sigma\subset\mathbb{R}^4$
\begin{enumerate}
\item[(i)] complete,
\item[(ii)] $J$-holomorphic
\item[(iii)] $m|_{T\Sigma}\geq 0$
\end{enumerate}
then $\Sigma$ is a plane. (Morally, the graph
of a linear function.)
\end{theorem}

\begin{theorem}[dim 3,4 Calabi, dim 5$\geq $
Pgorelov]
\label{theorem-smith-calabi-pogorelov}
$f:\mathbb{R}^2\to\mathbb{R}$
\begin{enumerate}
\item[(i)] $f$ convex,
\item[(ii)] $\det\operatorname{Hess}(f)=1$
(Monge-Ampère)
\end{enumerate}
Then $f$ is quadratic.

\begin{proof}
Short, done in seminar.
\end{proof}
\end{theorem}

And then we want to prove a compactness
property. So we take a sequence of surfaces
$(\Sigma_m,p_n)\subset S^1X$ with $p_n \in
\Sigma_n$. Suppose $\|\operatorname{I
I}_m\|\xrightarrow{m \to \infty}\infty$. Then
there exists $B_m$ and $q_m$ such that
\begin{enumerate}
\item[(i)] $\| \operatorname{ I
I}_m(q_m)\|=B_m$.
\item[(ii)] $B_m\xrightarrow{m \to
\infty}\infty$ and,
\item[(iii)] by a lemma that is easy to
prove, $\forall r\in
B_{\frac{1}{2\sqrt{\|\operatorname{I
I}_m(q_m)\|} }}(q_m)$.
\end{enumerate}
And then we rescale the metrics $g \to
g_m=B^2_mg$. That makes the metric be flatter
and flatter, like zooming in, and also makes
the shape operator of every surface have norm
1. and we get
$$
(S^1X,g_m)\longrightarrow(\mathbb{R}^5,)
$$
$$
\Sigma_m\longrightarrow \Sigma_m
\subset\mathbb{R}^5
$$
with respect to the famous Cheeger Gromov
topology in the limit, which is not so easily
defined. And by Arzelá-Ascoli and Elliptic
regularity magic (see M. Joshi Course notes)
(regularity is not smoothness but some
differentiability) the limit surface is
compact, positive, $J$-homolomorphic. And by
Jurgen's theorem they are flat. But that's a
contradiction with the fact that the shape
operators of these surfaces have norm 1.

So you have compactness. Yaaaaay!

\medskip\noindent
{\bf Editorial note.} Repairs made in
porting: in the horizontal and vertical
splitting, ``$TxX$'' was read as $T_xX$ and
the underbraces ``hor.'' and ``vert.'' were
moved to a sentence; the Saski metric
formula, which had an unclosed parenthesis
and no comma, was completed; the quadratic
form, written with a calligraphic lowercase
letter, is written $m$. The row of the
table with $0<k<1$ and the item
``$m((\xi,\mu),(\xi',\mu'))$'' (no right
hand side) are as in the original notes.
The numbered conditions use explicit labels.

\section{A primer on symplectic groupoids}
\label{section-prime-sympl-group}
%% speaker: Camilo Angulo, Jilin University
%% place: Geometric Structures Seminar, IMPA
%% date: February 13, 2025

\noindent
Camilo Angulo, Jilin University.
Geometric Structures Seminar, IMPA.
February 13, 2025.

\medskip\noindent
{\bf Abstract.}
In the late 17th century, Simeon Denis
Poisson discovered an operation that
helped encoding and producing conserved
quantities. This operation is what we
now know as a Lie bracket, an infinitesimal
symmetry, but what is its global
counterpart? Symplectic groupoids are one
possible answer to this question. In
this talk, we will introduce all the basic
concepts to define symplectic
groupoids, and their role in Poisson
geometry. We will discuss key examples, and
applications. The talk will be accessible to
those familiar with differential
geometry, but no prior knowledge of groupoids
will be assumed.


\medskip\noindent
{\bf Part 1. Poisson geometry.}

Hamiltonian formalism. Recall that being a
conserved quantity \(f \in
C^\infty(X)\) is the same thing as
\(\{H,f\}=0\).

\begin{itemize} \item We have seen that it is
always possible to take quotient
of a symplectic manifold with a group action
to obtain a  Poisson
manifold.  \item Then we have found a way to
produce a symplectic foliation
from a 2-vector \(\pi \in
\mathfrak{X}^2(M):=\Lambda^{2}(TM)\).  \item

\begin{remark}
\[\left\{ \text{Lie algebra on
\(\mathfrak{g}\)}
\right\} \xrightarrow{1-1}\left\{
\text{Linear Poisson bracket on
\(\mathfrak{g}^*\)}  \right\}  \]
\end{remark}

\item We saw very nice examples
of foliation that have to do with Lie
algebras. So \(\mathfrak{b}^*_3\) which
gives the ``open book foliation", and
\(\mathfrak{e}^*\) that gives a foliation
by cylinders.
\end{itemize}

\medskip\noindent
{\bf Part 2. Symplectic realizations.}

Consider
\[(\Sigma,\omega) \xrightarrow{\mu}(M,\pi)\]


\iffalse\begin{tikzcd}
	T_p\Sigma
	\arrow[r,"d_p\mu"]\arrow[d,"\omega^\flat"]&
	T_xM\\
	T_p^*\Sigma &
	T^*_pM\arrow[l,"(d\mu)^*",swap]\arrow[u,"\pi
	^\sharp"]
\end{tikzcd}\]\fi
So that
\[\pi ^\sharp=d_p\mu \circ \omega^{-1}
\circ(d\mu)^*\]



\begin{lemma}
\(\dim (\Sigma) \geq  2 \dim (M) -
\operatorname{rk}(\pi_x)\) for all \(x \in
M\).
\end{lemma}

\begin{proof}
Done in seminar.
\end{proof}

\begin{example}
\((\mathbb{R}^2,0)\). So the map
$$
\begin{aligned}
	(\mathbb{R}^4,dx \wedge du + dy \wedge dv)
	&\longrightarrow  \mathbb{R}^2\\
	(x,y,u,v) &\longmapsto (x,y)
\end{aligned}
$$

\end{example}

\begin{exercise}
Find the symplectic realization \(\omega\) in
\((\mathbb{R}^4, \omega )
\xrightarrow{\mu}(\mathbb{R}^2,(x^2+y^2)
\partial_x
\wedge \partial_y\)

$$
\begin{aligned}
	(\mathbb{R}^4, \omega) &\longrightarrow
	(\mathbb{R}^2,(x^2+y^2) \partial_x \wedge
	\partial _y) \\
	(x,y,z,w) &\longmapsto (x,y)
\end{aligned}
$$

Also find the symplectic realization of
\(\operatorname{ a f f }^*\) with
\(\{x,y\}=x\).
\end{exercise}

\medskip\noindent
{\bf Part 3: Groupoids}

\medskip\noindent
{\bf Motivation}
\begin{enumerate}
\item Fundamental grupoid: objects are points
in the manifold and arrows are paths.
\item The action of \(S^1\) on \(S^2\) by
rotation does not give a nice quotient
because there are two singular points.
Consider the groupoid  \(S^1 \times S^2\) of
orbits. These are the arrows. The points are
just the points of \(S^2\).
\item Consider a foliation (like Möbius
foliation of circles; where there is a
singular circle, the soul). You can do the
same thing as in fundamental groupid
leafwise. Arrows then are equivalence classes
of paths that live inside leaves. This is
called monodromy of a foliation. Again
objects are points.
\item You can take a quotient of monodromy
using a connection given by the foliation.
This allows to identify certain paths between
the leaves. This is called {\it holonomy (of
a foliatiation)}. (So you can make this
notion match the usual holonomy given by
riemmanian connection.)

{\bf Upshot.}
	So the point is taking some sort of function
	space on these groupoids you can gather the
	information given by the non-smooth quotient
	(like in the case of the sphere rotating).
	So this groupoid motivation says how to get
	some structure that resembles a non smooth
	quotient.
\item Last motivation: the grid of squares
has a tone of symmetries. If you restrict to
just a few squares you loose so many
symmetries. But there's a grupoid hidden in
there that tells you what you intuition knows
about this finite grid of squares.
\end{enumerate}

\begin{definition}
A {\it groupoid} is a category where all
morphisms are invertible.
\end{definition}

So there is a kind of product among the
objects, given by composition but:
 not every two pair of objects can be
 multiplied!---only those whose
source and target match. So that's the lance
about grupoids.

Just so you make sure you understand: the
groupoid \(G\) is the morphisms of the
category. The objects are points (of a
manifold).

\begin{definition}
{\it Lie grupoid} is when the following
diagram is inside category of smooth
manifolds and \(s,t\) (source and target
maps) are submersions:
$$
\xymatrix{
G^{(2)}\ar[r]^m& G\ar[r]^ s_t & M\ar[r]^u&G
}
$$
\end{definition}

\begin{proposition}[Properties of Lie
groupoids]
\begin{itemize}
\item \(m\) is also a submersion.
\item \(i\) (inversion) is a diffeomorphism.
\item  \(u\) (unit=identity) is an embbeding.
\end{itemize}
\end{proposition}

\begin{definition}
Consider \(x \in M\) and the inverse image of
source map: \(s^{-1}(x)=\{\text{arrows that
start at \(x\)} \}\). Now if you act with
\(t\) on this set you get {\it the orbit} of
\(x\): \(\{\text{\(y \in M\) such that there
is an arrow from \(x\) to \(y\)} \}\).

And there also {\it an isotropy} \(G_x=\{g
\in G: \text{\(g\) goes from \(x\) to \(x\)}
\}\)
\end{definition}

\begin{example}
\begin{enumerate}
\item \(G=M\),  \(M=M\).
 \item Lie groups.
	\item Lie group bundles.
\item \(G=M \times M\), \(M =M\).
 \item Fundamental groupoid. Isotropy group
 is fundamental group! And orbit
 is…\clearpage universal cover!
\item Subgroupoids.
\item  Foliations.
\item If you have a normal group action of
\(G\) on \(M\) you construct a groupoid
action with groupoid \(G \times M\) and
objects \(M\), with product given on the
group part of the product. Orbits are orbits.
Isotropy group is isotropy group.
 \item Principal bundles.
\end{enumerate}
\end{example}

{\bf Back to Poisson.}

There's also a notion of Lie algebroid. Which
is strange. But the point is that to every
Poisson manifold there is a Lie algebroid.

So the question is whether there is a Lie
groupoid associated to that Lie algebroid.
Not always.

\medskip\noindent
{\bf Big question}[Fernandez and ?]
When a symplectic manifold is integrable?


(Remember that integrating means go from
algebra(oid) to group(oid).

And the point is that:

\medskip\noindent
When you {\it can} go back, you get a
\textit{ symplectic groupoid}.

\begin{remark}
Look for Kontsevich's notes on Weinstein!
\end{remark}

\begin{remark}
History: Weinstein did this intending to do
quantization (geometric?) on Poisson
manifolds. (That involves a \(C^*\)algebra
coming from the symplectic groupoid.)
\end{remark}

\begin{definition}
A {\it symplectic groupoid} is a groupoid
\(G, M\) together with  \(\omega \in
\Omega^2(M)\) such that \(\omega\) is
symplectic and multiplicative, meaning that
\(\partial  \omega =0\), that is, \(\iff
m^*\omega=\operatorname{pr}_1^*\omega+
\operatorname{pr}^*_2\omega\in
\Omega^{2}(G^{(2)})\iff\) take two vectors \(
X_k,Y_k \in TG\), and \(\omega(X_0 \star
Y_0,X_1\star Y_1) =  \omega(X_0,Y_0) + \omega
(X_1,Y_1)\)
\end{definition}

\begin{theorem}
If \((G, \omega)\) is a symplectic groupoid,
then
\begin{enumerate}
\item \(\exists !\) poisson structure on
\(M\)
\item for which \(t: G \to M\) is a
symplectic realization,
\item  Leaves are connected components of
orbits,
\item \(\mathsf{Lie}(G)\cong T^*M\) via \(X
\mapsto -u ^*(i_X \omega\).
\end{enumerate}
\end{theorem}

\begin{remark}
Look for Alejandro Cabrera, Kontsevich. There
are two things one is de Rham and the other…
from the future: simplicial?
\end{remark}

\medskip\noindent
{\bf Upshot.} The obstruction to knowing when
symplectic groupoid exists is ``variation of
symplectic form \(\omega=(1+t^2)
\omega_{S^2}\)". So how does the symplectic
group vary from leaf to leaf. So there are
two situations in which the thing doesn't
work.

\section{Geometria de Poisson: primeiras duas
aulas}
\label{section-geometria-de-poisson}
%% speaker: Henrique Bursztyn, IMPA
%% place: IMPA course
%% date: March 11 and 13, 2025

\noindent
Henrique Bursztyn, IMPA.
IMPA course.
March 11, 13, 2025.

\medskip
Livros:
\begin{itemize}
\item Lectures em PG, Carnic-Gernandes-Marcit
\item A brief introduction to PG, HB
\end{itemize}

\medskip\noindent
{\bf História.}
Em 1809, Poisson buscava uma formulação
geométrica da mecânica clássica (celeste).
Espaço fase, função Hamiltoniana, campo
Hamiltoniano, equações de Hamilton. O
primeiro colchete de Poisson é
$$
\begin{aligned}
\{\cdot ,\cdot \}:
C^\infty(\mathbb{R}^{2n})\times
C^\infty(\mathbb{R}^{2n}) &\longrightarrow
C^\infty(\mathbb{R}^{2n}) \\
\{f,g\} &=\sum_{i=1}^n \frac{\partial
f}{\partial p_i}\frac{\partial g}{\partial
q_i}-\frac{\partial f}{\partial
q_i}\frac{\partial g}{\partial p_i}
\end{aligned}
$$

Interpretação dinâmica:
$\{H,f\}=\mathcal{L}_{X_H}f$, que como temos
antisimetria, é o mesmo que
$-\mathcal{L}_{X_f}H$.

Conhecer a dinâmica do sistema é resolver às
equações de Hamilton, EDOs. Integrais
primeiras. Aí nasce a área de \emph{sistemas
completamente integraveis}.

\begin{theorem}[de Poisson, 1809]
\label{theorem-bursztyn-poisson}
O colchete de Poisson de duas integrais
primeiras é uma integral primeira, i.e.
$$
\{H,f\}=0,\{H,g\}=0 \implies \left\{
\{f,g\},H \right\} =0
$$
\end{theorem}

Mas foi o Jacobi que descobreu o meolho
daquele teorema:
\begin{theorem}[Jacobi, 1842]
\label{theorem-bursztyn-jacobi}
$$
\{h,\{f,g\}\} +\{g,\{h,f\}\} +\{f,\{g,h\}\}
=0
$$
\end{theorem}
De fato, o teorema de Poisson segue do
teorema de Jacobi.

Em 1880 S. Lie trabalha em \emph{álgebras de
Lie}. Ele mostra que toda álgebra de Lie vem
de um grupo de Lie localmente. Isso levou à
\emph{teoria de Lie} onde foram definidas as
estruturas de Poisson.

Em 1970, Kirillov, Sourieau e Kostant voltam
a trabalhar no colchete de Poisson. Dois
artigos que marcaram a era moderna da
geometria de Poisson são Lichnerowicz (?) e
Weinstein (1983).

\medskip\noindent
{\bf Motivação para a geometria de Poisson.}
\begin{itemize}
\item Mecânica geometrica (plasmas)/Teoria de
campos.
\item Sistemas integráveis $\to$ sistemas
bihamiltonianos Poisson-Nijenhuis (algum
scenário onde temos duas estruturas de
Poisson compatíveis\ldots{}).
\item Teoria de representações/grupos
quânticos (Drinfeld, Fadeev)
$\rightsquigarrow$ grupos de Lie/Poisson.
\item Quantização por deformação. Passagem do
formalismo clássico com o colchete de
Poisson, para um formalismo com uma álgebra
não comutativa $\mathbb{A},*$.
\end{itemize}

\medskip\noindent
{\bf Definições.}

\begin{definition}
\label{definition-bursztyn-poisson-bracket}
$M$ variedade diferenciável. Um
\emph{colchete de Poisson} em $M$ é uma
operação $\mathbb{R}$-bilinear
$$
\{\cdot ,\cdot \}: C^\infty(M)\times
C^\infty(M) \longrightarrow C^\infty(M)
$$
tal que
\begin{enumerate}
\item (Antisimetria.) $\{f,g\}=-\{g,f\}$.
\item (Jacobi.) $\{h,\{f,g\}\} +\{g,\{h,f\}\}
+\{f,\{g,h\}\} =0$.
\item (Leibniz.)
$\{f,gh\}=\{f,g\}h+\{f,h\}g$.
\end{enumerate}
\end{definition}
\begin{remark}
\label{remark-bursztyn-lie-algebra}
As condicões (1) e (2) dizem que
$(C^\infty(M),\{\})$ é uma álgebra de Lie.
Como ainda é uma álgebra comutativa com o
produto usual, a condição (3) diz como é que
interagem esses dois produtos.
\end{remark}
\begin{remark}
\label{remark-bursztyn-poisson-algebra}
Para qualquer álgebra comutativa
$\mathcal{A}$ podemos introduzir um colchete
de Lie e pedir a condição Leibniz, e isso se
chama de \emph{álgebra de Poisson}. E se o
colchete satisfaz (1) e (3), se chama de
\emph{estrutura quase Poisson}, que já tem um
significado geométrico.
\end{remark}
\begin{definition}
\label{definition-bursztyn-poisson-map}
\emph{Aplicação (ou morfismo) de Poisson}
entre $(M_1,\{\cdot ,\cdot \}_1)$ e $(M_2,\{
\cdot ,\cdot \}_2)$ é $\varphi:M_1 \to M_2$
que preserva colchetes, i.e. o pullback
$$
\begin{aligned}
\varphi^*: C^\infty(M_2) &\longrightarrow
C^\infty(M_1) \\
f &\longmapsto f \circ \varphi
\end{aligned}
$$
preserva colchete no sentido de que
$$
\{f,g\}_2 \circ \varphi=\{f \circ \varphi,g
\circ \varphi\}_1.
$$
\end{definition}
Temos um campo que chamamos de
\emph{hamiltoniano} que podemos tirar da
condição leibniz do colchete. Isso é porque
$\{f,\cdot \}:C^\infty(M) \to C^\infty(M)$ é
uma derivação! Definindo esse campo como
$X_f$ obtemos
$$
\{f,g\}=\mathcal{L}_{X_f}f=dg(X_f)=-df(X_g)
$$

\begin{remark}
\label{remark-bursztyn-local}
Os colchetes são locais no sentido de que
podemos restringir num aberto.
\end{remark}
\begin{remark}
\label{remark-bursztyn-casimir-zero}
$\{f,f\}=0$ ($f$ é preservada por seu campo
hamiltoniano).
\end{remark}

\medskip\noindent
{\bf Terminologia.}
\begin{itemize}
\item $f,g \in C^\infty(M)$ estão em
\emph{involução} se $\{f,g\}=0$.
\item Se $f$ é tal que $\{f,g\}=0 \forall g
\iff X_f=0$, então $f$ é dita de
\emph{Casimir}. No caso simplético isso não
faz muito sentido---todas as funções são
Casimir porque a forma simplética é não
degenerada. Mas no caso geral isso pode
mudar.
\end{itemize}

\begin{example}
\label{example-bursztyn-poisson-brackets}
\begin{itemize}
\item $(M,\{\cdot ,\cdot \}\equiv0)$.

Outro de jeito de definir funções
$f_1,\ldots,f_k$ com $\{f_i,f_j\}=0$ é
simplesmente dizer $(M,\{\cdot ,\cdot
\})\xrightarrow{F}(\mathbb{R}^k,\{\cdot
,\cdot \} \equiv 0)$ é uma aplicação de
Poisson.
\item $\mathbb{R}^{2n}=\{(q,p)\}$, o colchete
de Poisson de Poisson é um colchete de
Poisson.
\item Mais geralmente, toda vez que tenha uma
variedade simmplética $(M,\omega)$, então
temos o campo hamiltoniano, e podemos definir
como você já sabe
$\{f,g\}:=\omega(X_g,X_f)=dg(X_f)=-df(X_g)$.
É isso é um colchete de Poisson, para vê-lo
vai notar que Jacobi é equivalente a
$d\omega=0$.

Qualquer colchete de Poisson que vem de uma
estrutura simplética pode ser escrito como o
colchete de Poisson de Poisson em coordenadas
locais pelo teorema de Darboux.

\item Considere o produto de $(M_1,\{\cdot
,\cdot \}_1)$ e $(M_2,\{ \cdot ,\cdot \}_2)$
é $\varphi:M_1 \to M_2$, $M_1 \times M_2$.
Então
$$
\{f,g\}(x_1,x_2)=\{f(\cdot ,x_2),g(\cdot
,x_2)\}_1(x_1)+\{f(x_1,\cdot ),g(x_1,\cdot
)\}_2(x_2).
$$
{\bf Exercício!} Mostre que esse é um
colchete de Poisson. Ainda, que as projeções
são mapas de Poisson e os pullbacks de
funções em cada $M_i$ Poisson-comutam no
produto, i.e.
$$
\{p_1^*C^\infty(M_1),p_2^*C^\infty(M_2)\}=0
$$
[Figure in the original notes: commutative
triangle with $M_1 \times M_2$ on top,
mapping to $M_1$ by $p_1$ and to $M_2$ by
$p_2$.]
\item $S$ variedade, considere uma família
suave de estruturas simpléticas, i.e.
$\omega_t \in \Omega^2(S)$ simpléticas.
Considere as estruturas de Poisson associadas
a cada uma delas, i.e. $\{\cdot ,\cdot \}_t$.
Então pode produzir um colchete em $M:=S
\times \mathbb{R}$ dado por
$$
\{f,g\}(x,t):=\{f(\cdot ,t),g(\cdot
,t)\}_t(x).
$$
Creo que es sólo pegar la estructura de
Poisson para cada $t$.
\end{itemize}
\end{example}

\medskip\noindent
{\bf Ponto de vista tensorial.}
Para isso precisamos de \emph{multivetores},
que é o análogo das formas diferenciais para
campos vetoriais, i.e. seções do fibrado
$\Lambda^{\bullet}(TM)$. I.e.,
$\mathfrak{X}^\bullet(M)=
\Gamma(\Lambda^{\bullet}(TM))$.

Considere uma estrutura $\{\cdot ,\cdot \}$
quase-Poisson em $M$. Então existe um único
bivetor $\pi \in \Gamma(\Lambda^{2}(TM))$ tal
que
$$
\{f,g\}=\pi(df,dg)
$$
O lance é que o colchete quase-Poisson não
depende das funções, mas sim das diferenciais
delas.

De fato, temos uma correspondência entre
colchetes quase-Poisson $\{\cdot ,\cdot \}$ e
bivetores $\pi \in \Gamma(\Lambda^{2}(TM))$
em $M$:
$$
\{\cdot ,\cdot \} \leftrightsquigarrow \pi
$$
Dizemos que $\pi \in \Gamma(\Lambda^{2}(TM))$
é um \emph{bivetor de Poisson} se
$\{f,g\}=\pi(df,dg)$ satisfaz Jacobi. Em
coordenadas locais $(x_1,\ldots,x_n)$, um
bivetor $\pi$ sempre pode ser escrito como
$$
\begin{aligned}
\pi &=\frac{1}{2}\sum_{i,j}\pi_{i,j}(x)
\frac{\partial }{\partial x_i}\wedge
\frac{\partial }{\partial x_j} \\
&=\sum_{i<j}\pi_{ij}(x)\frac{\partial
}{\partial x_i}\wedge\frac{\partial
}{\partial x_j}
\end{aligned}
$$
de modo que qualquer colchete quase-Poisson,
localmente se escreve desse jeito:
$$
\{f,g\}(x)=\sum_{i,j}\pi_{ij}(x)
\frac{\partial
f}{\partial x_i}\frac{\partial g}{\partial
x_j}
$$
onde (creo que), em coordenadas locais,
$\pi_{ij}(x)=\{x_i,x_j\}$.
\begin{remark}
\label{remark-bursztyn-jacobi-pde}
A condição de ser Jacobi nas funções
$\pi_{ij}$ é uma EDP muito difícil, tem
outras formas de mostrar que algo é Jacobi
além de resolvê-la.
\end{remark}

\medskip\noindent
{\bf Editorial note.} Repairs made in
porting: in the example list, the last item
(family of symplectic structures) stood
after the closing of the list inside the
example; the list now includes it. The
correspondence between quasi-Poisson
brackets and bivectors, written in the
original with words inside a formula, is
now a sentence and a formula. The formula
for the Hamiltonian vector field, as in the
original, has $\mathcal{L}_{X_f}f$ where
$\mathcal{L}_{X_f}g$ is presumably meant; it
is left as written.

\section{Lie Groupoids: Foundations,
Advances,
and Future Directions}
\label{section-lie-groupoids-foundations}
%% speaker: Matias del Hoyo, UFRJ
%% place: Seminário das Sextas
%% date: March 21, 2025

\noindent
Matias del Hoyo, UFRJ.
Seminário das Sextas.
March 21, 2025.

\medskip\noindent
{\bf Abstract.}
Lie groupoids are central to modern geometry
and mathematical physics, connecting
differential geometry, topology, and
dynamical systems. Based on a recent chapter
co-authored with H. Bursztyn for the
\href{https://www.sciencedirect.com/science/article/abs/pii/B9780323957038000240}{Encyclopedia
of Mathematical Physics} this talk will
explore the foundations of Lie groupoids,
their interactions with related structures
like Lie algebroids and Poisson manifolds,
and highlight recent advances and emerging
research directions in the field.

\medskip\noindent
{\bf Lie groupoids.}

\begin{definition}
\label{definition-hoyo-lie-groupoid}
A \emph{Lie groupoid} $G \rightrightarrows M$
consists of manifolds $G$, $M$, submersions
$s,t:G \to M$ and a multiplication $m$ with
unit $u$ and inverse $i$
$$
\begin{aligned}
m: G \times_M G &\to G \qquad (z
\overset{g}{\leftarrow }y,
y\overset{g}{\leftarrow } x )\mapsto
(z\overset{hg}{\leftarrow }x) \\
u:M &\to G \qquad \qquad x \mapsto (x
\overset{1_x}{\leftarrow }x) \\
i:G &\to G \qquad \qquad (y
\overset{g}{\leftarrow }x)\mapsto
(x\overset{g^{-1}}{\leftarrow }y)
\end{aligned}
$$
A \emph{morphism}
$\phi:(\tilde{G}\rightrightarrows M) \to (G
\rightrightarrows M)$ consists of maps
$\phi_1:\tilde{G} \to G$, $\phi_2:\tilde{M}
\to M$ preserving source, target,
multiplication and units.
\end{definition}

\medskip\noindent
{\bf Examples.}

\begin{example}
\label{example-hoyo-groupoids}
\begin{itemize}
\item For $M$ manifold:
\begin{itemize}
\item $M\rightrightarrows M$ unit groupoid.
\item $M \times M \rightrightarrows M$ pair
groupoid.
\item $\pi_{1}(M)\rightrightarrows M$
fundamental groupoid.
\end{itemize}
\item $G$ Lie group: $G \rightrightarrows *$
Lie group

\item $q:M \to N$ submersion: $M \times_M
\rightrightarrows M$ submersion groupoid. Ex:
$\mathcal{U}=\{U_i\}$ open cover
$\rightsquigarrow$ $\coprod
U_{ji}\rightrightarrows\coprod U_i$ \v Chech
groupoid.
\item $G \curvearrowright M$ action: $G
\times M \rightrightarrows M$ action
groupoid.
\item $X \in \Gamma(TM)$: $U_X \subset
\mathbb{R}\times M \rightrightarrows M$ flow
gruopoid.
\item $P \to M$ principal bundle: $(P\times
P)/G \rightrightarrows M$ gauge groupoid.
\item $F \subset M$ foliation:
\begin{itemize}
\item $\operatorname{Mon}F \rightrightarrows
M$ monodromy groupoid
\item $\operatorname{Hol}M \rightrightarrows
M$ holonomy groupoid.
\item $V \to M$ vector bundle:
$\mathsf{GL}(V)\rightrightarrows M$ bigeneral
linear groupoid.
\end{itemize}
\end{itemize}
\end{example}

\medskip\noindent
{\bf Actions and representations.}
Here's one of the key ideas: relation between
group and groupoid. Groups in geometry serve
to model symmetries of a space: a group $G
\rightrightarrows *$ $\leftrightsquigarrow$
symmetries of a space $E \to *$. And
groupoids model symmetries of a \emph{family}
parametrized by $M$: a groupoid $G
\rightrightarrows M$ $\leftrightsquigarrow$
symmetries of a family $E \to M$.

\begin{definition}
\label{definition-hoyo-groupoid-action}
Given $G\rightrightarrows M$ a Lie grouoid
and $E \to M$ a surjective submersion, a
\emph{groupoid action} $\rho:G
\curvearrowright E$ is a map such that
$\rho_{1_x}=\operatorname{id}_{E_x}$ and
$\rho_h\rho_g=\rho_{gh}$:
$$
\rho:G \times_M E \to E \qquad (y\leftarrow
x,e/x)\mapsto \rho_g(e)/y
$$
It is a \emph{representation} if $E \to M$
vector bundle and $\rho_g E_x \to E_y$
linear.
\end{definition}

\begin{example}
\label{example-hoyo-actions}
\begin{itemize}
\item The \emph{parallel action}. Take a
vector bundle with a flat connection. Then
the fundamental groupoid \emph{acts by
parallel transport}, i.e. $(E \to M, \nabla)$
is a groupoid action
$(\pi_{1}(M)\rightrightarrows M
\curvearrowright (E \to M)$. Notice this is
\emph{not} a group action, it's naturally a
groupoid action.
\item A \emph{Hamiltonian action} $G
\curvearrowright(M,\omega)$ on a symplectic
manifold is an action of the action groupoid
$(G \times
\mathfrak{g}^*\rightrightarrows\mathfrak{g}^*
\curvearrowright(M\to
\mathfrak{g}^*)$.
\item The \emph{(linear) holonomy} of a
foliation $F \subset TM$ is a representation
of the monodromy groupoid
$(\operatorname{Mon}(F) \rightrightarrows M
\curvearrowright(TM/F \to M)$.
\end{itemize}
\end{example}

\medskip\noindent
{\bf Groupoid fibrations.}
A representation $(G\rightrightarrows
M)\curvearrowright(E \to M)$ gives rise to a
groupoid morphism with a (unique) lift of
arrows:

[Figure in the original notes: commutative
diagram with $G \times_M E \rightrightarrows
E$ mapping down to $G\rightrightarrows M$; an
arrow $x \to y$ labelled $g$ in $M$ lifts to
the arrow $e \to \rho(g,e)$ in $E$ (in red).]

A \emph{fibration} $(\Gamma \rightrightarrows
E) \to (G \rightrightarrows M)$ is a morphism
with (non-necessarily unique) lift of arrows.
A \emph{VB-groupoid} is a linear fibration
(called like this for historical reasons).

\begin{example}
\label{example-hoyo-fibrations}
\begin{itemize}
\item A usual representation $(G \ltimes_E
\rightrightarrows E) \to (G \rightrightarrows
M)$
\item The tangent and the cotangent groupoid
$(TG \rightrightarrows TM) \to (G
\rightrightarrows M)$ and $(T^*G
\rightrightarrows A^*) \to
(G\rightrightarrows M)$
\end{itemize}
\end{example}
\begin{theorem}[Grothendieck;Gacria-Saz,
Mehta;dH,Ortiz]
\label{theorem-hoyo-vb-groupoids}
There is a 1-1 correspondence between
VB-groupoids $\Gamma\to E$ over $G
\rightrightarrows M$ and
\emph{representations up to homotopy}
$G\curvearrowright(C \oplus E \to M)$.
\end{theorem}

\medskip\noindent
{\bf Nerve and cohomology.}
Given $G \rightrightarrows M$ Lie groupoid,
its \emph{nerve} $(NG_n,d_i,s_j)$ is the
simplicial manifold
\begin{itemize}
\item $n$-simplices: $NG_n=\{x_n
\overset{g_n}{\leftarrow}x_{n-1}
\overset{g_{n-1}}{\leftarrow}\cdots
\overset{g_2}{\leftarrow}x_1
\overset{g_1}{\leftarrow}x_0\}$
\item Faces $d_i:NG_n \to NG_{n-1}$ compose
two arrows (or forget one)
\item Degeneracies $s_j:NG_n \to NG_{n+1}$
insert and identity.
\end{itemize}
The \emph{groupoid DGA} is
$C(G)=(\operatorname{sin}(NG_\bullet),\cup
,\partial=\sum_{i}(-1)^i d_i^*)$. The
\emph{groupoid cohomology} is
$H^{n}(G)=H^{n}(C(G))$.

\medskip\noindent
{\bf Linearization problem.}
\begin{remark}
\label{remark-hoyo-proper}
Here was discussed a notion of \emph{proper
groupoid}, which have analogous properties of
proper Lie groups: quotient is Hausdorff and
admits a Haar measure.
\end{remark}

\medskip\noindent
{\bf Metric on a Lie groupoid.}

\begin{definition}
\label{definition-hoyo-metric}
A \emph{metric} on $G \rightrightarrows M$ is
a metric $\eta^{(2)}$ on $NG_2=G \times_M G$
such that the face maps are Riemannian
submersions\ldots{}
\end{definition}
It's actually a \emph{system of metric},
making source, target, multiplication
Riemannian submersions.

\begin{theorem}[dH-Fernandes, 2018]
\label{theorem-hoyo-metric}
Every proper Lie groupoid admits a metric,
and a metric yields linearizations through
the exponential maps.
\end{theorem}

\begin{remark}
\label{remark-hoyo-invariant-metric}
How to construct a invariant metric from any
metric.
\end{remark}

\begin{remark}[About the future]
\label{remark-hoyo-future}
It's not yet clear how to define a proper
higher Lie groupoid. Speaker is trying to
develop geometry for higher Lie groupoids.
\end{remark}

\medskip\noindent
{\bf Part 2.}

\medskip\noindent
{\bf Abstract Lie algebroids.}

\begin{definition}
\label{definition-hoyo-lie-algebroid}
A \emph{Lie algebroid} $A \implies M$ is a
vector bundle $A \to M$ with a Lie bracket
$[\cdot,\cdot]$ on $\Gamma(A)$ and a linear
map $\rho:A \to TM$ satisfying Leibniz rule
$$
[\alpha,f
\beta]=f[\alpha,\beta]+\rho(\alpha)(f)\beta
\qquad \qquad \alpha, \beta \in \Gamma(A),f
\in C^\infty(M)
$$
If
\end{definition}

\begin{remark}
\label{remark-hoyo-higher-algebras}
Also \emph{higher Lie algebras} have been
defined.
\end{remark}

\medskip\noindent
{\bf Last slide.}
There were several slides in between\ldots{}

\begin{theorem}[Crainic-Fernandes]
\label{theorem-hoyo-crainic-fernandes}
There is a 1-1 correspodnence between
($s$-simply connected) symplectic Lie
groupoids and (integrable) Poisson manifolds.
\end{theorem}

Why symplectic groupoid and Poisson
manifolds? They seem to be different objects.
Theese (which?) induce Lie algebra. Because
cotangent structure

\begin{quotation}
``Poisson manifold induce a symplectic lie
algebra.'' ``Integration of symplectic
algebra is symplectic group''
\end{quotation}

\medskip\noindent
{\bf Editorial note.} In porting, the
arrows written with two stacked ``to''
arrows are written $\rightrightarrows$,
and the group action symbol (written with a
blackboard letter ``y'' in the original) is
written $\curvearrowright$. The original
notes break the pieces ``Part 2'' and
``Last slide'' off as a subsection and a
section; here they are headings of this
section. The definition of Lie algebroid
ends in the dangling word ``If'', and the
$M \times_M$ in the submersion groupoid
is as in the original notes; neither is
repaired. The diagram of a fibration is
replaced by a bracketed description.

\section{On wrapped Floer homology barcode
entropy and hyperbolic sets}
\label{section-wrapp-floer-homol-barco}
%% speaker: Rafael Fernandes, UC Santa Cruz
%% place: Differential Geometry Seminar, IMPA
%% date: September 4, 2025

\noindent
Rafael Fernandes, UC Santa Cruz.
Differential Geometry Seminar, IMPA.
September 4, 2025.

\medskip\noindent
{\bf Abstract.} In this talk, we will discuss
the interplay between the
wrapped Floer homology barcode and
topological entropy. The concept of barcode
entropy was introduced by Çineli, Ginzburg,
and Gürel and has been shown to be
related to the topological entropy of the
underlying dynamical system in various
settings. Specifically, we will explore how,
in the presence of a topologically
transitive, locally maximal hyperbolic set
for the Reeb flow on the boundary of
a Liouville domain, barcode entropy is
bounded below by the topological entropy
restricted to the hyperbolic set.

\medskip\noindent

Let $M^n$ be a manifold.
$\omega \in \Omega^2(M)$ is {\it symplectic}
if $d \omega=0$ and it is
nondegenerate.

\begin{example}
\label{example-Rn-is-symplectic}
$\mathbb{R}^n$ is symplectic with canonical
Darboux form.
\end{example}

Recall the definition of Hamiltonian vector
field
associated to a function $H \in C^\infty(M)$.

\begin{definition}
\label{definition-nonde-diffe}
A diffeomorphism $\varphi:M \to M$ is called
{\it non-degenerate} if
$\Phi(\varphi)\cap\Delta \subset M\times M$
(pitchfork, i.e. transversal intersection!).
\end{definition}

Let $M^{2n}$ be a closed symplectic manifold.
Arnold's conjecture says
\begin{enumerate}
\item If $\varphi=\varphi_H$ (Hamiltonian
flow) is nondegenerate, then
$$
\# \text{Fix}(\varphi_H) \geq
\sum_{i=0}^{2n}\dim H_i(M,k)=\dim H_*(M,k)
$$
\item If $\varphi=\varphi_H$ is degenerate,
then
$$
\# \text{Fix}(\varphi)\geq \text{Cl}(M)+1
$$
where $\text{Cl}(M)$ is the maximum number of
homology classes
we can add before getting to zero.
\end{enumerate}

Why do we care? Because
$$
\#
\text{Fix}(\varphi_H)
\leftrightarrow\{\text{1-periodic
orbits of $X_H$}\}
$$

Idea by Floer. Construct an invariant that
would
say something about periodic orbits.

\medskip\noindent
{\bf Question.} Can Floer theory capture
other ``dynamical information''?
(Other than the periodic orbits.)

\medskip\noindent
A {\it persistence module} is a pair
$(V,\Pi)$, where $V=\{V_t\}_{t \in
\mathbb{R}}$ is a family of
$\mathbb{F}$-vector spaces
and $\Pi=\{\Pi_{st}\}_{s \leq  t}$ is a
family of maps such that
\begin{enumerate}
\item $\Pi_{ss}=, \Pi_{ts}\circ
\Pi_{rt}=\Pi_{rs}$.
\item $\exists  s \subset \mathbb{R}$ such
that
$\Pi_{s t}$ is an isomorphism for $s,t$ in
the same connected component of
$R\setminus S$.
\item $\Pi_{s t}$ have finite rank.
\item $\exists s_0$ $V_s=\{0\}$, $s \leq
s_0$.
\item $V_t= \lim_{s \to t} V_s$ (lower
limit!!)
\end{enumerate}

\begin{theorem}
\label{theorem-persi-modul-sum-integ}
Any persistence module is a sum of integral
persistence modules,
$$
(V,\Pi) \cong \bigoplus_{I \in B(V)}F(I).
$$
\end{theorem}

\begin{example}
\label{example-hart-and-sphere}
Heart and sphere. There is a noise in the
persistence module
of the heart due to an unnecessary critical
point.
\end{example}

$(M^{2n},\omega)$ a {\it Liouville domain} is
a compact
symplectic manifold and $X \in
\mathfrak{X}(M)$
with $X \cap\partial M$ (pitchfork, i.e.
transversal intersection!)
pointing outwards
and preserved by the symplectic form, i.e.
$\mathcal{L}_X \omega=\omega$
($\omega=d \alpha$).

When we restrict $\omega$ to the boundary,
we obtain a contact form and get some
interesting dynamics.

A Lagrangian $(L,\partial L) \subset
(M,\partial M)$
is {\it asympotically conical} if
\begin{enumerate}
\item $\partial L \subset \partial M$ is
Legendrian.
\item $L \cap [1-\varepsilon,1] \times
\partial M=
[1-\varepsilon,1] \times \partial L$.
\end{enumerate}

\begin{remark}
\label{remark-Hamiltonian}
Take a Hamiltonian $H: \hat{M}\to \mathbb{R}$
such that
$$
\begin{cases}
H(r,x)=h(r)\qquad &r=1 \\
H(r,x)=rT-B\qquad &
\end{cases}
$$
then $X_H=h'(r)R_\alpha$.

For $L_0,L,A\subset$ Lagrangians, $H$ linear
at infinite,
then $A_{H}^{L_0 \to L_1}$,
$$
\begin{aligned}
A_{H}^{L_0 \to L_1}: P_{L_0\to L_1}
&\longrightarrow \mathbb{R} \\
\gamma &\longmapsto \int_0^1 \gamma^* \alpha-
\int_0^1H(x(t))dt
\end{aligned}
$$
where $P_{L_0\to L} =\{\gamma:[0,1] \to
\hat{M}:\gamma(0) \in L_0,
\gamma(1) \in L_1\}
$ is the set of chords.
\end{remark}

\begin{remark}
\label{remark-1-chords-are-crticial-points}
$\text{crit}(A_{H}^{L_0 \to
L_1})=\{\text{1-chords of $X_H$ from
 $L_0$ to $L_1$}\}$.
\end{remark}

Putting a metric on $P_{L_0 \to L_1}$ we can
consider $\varphi:\mathbb{R} \times
[0,1] \to \hat{M}$, solutions of some PDE
which is some kind
of generalization of a gradient, $- \nabla
A_{H}^{L_0 \to L_1}$.
These solutions can be put in a moduli space
$$
\tilde{\mathcal{M}}(x_-,x_+,H,J)=
\{\varphi\text{
solutions s.t. …}\}
$$
Then we define a boundary operator
$\partial$.
\begin{theorem}
\label{theorem-homology}
$\partial^2=0$
\end{theorem}

So that we have a homology, called {\it
wrapped Floer homology}
$HW^t(H,L_0,L_1,J)$

\begin{remark}
\label{remark-embedding}
We have $H \leq  K \rightsquigarrow
HW^t(H,L_0,L_1,J) \to
HW^t(K,L_0,L_1,K)$.
\end{remark}

\begin{definition}
\label{definition-direct-limit}
For $t \geq 0$
$$
HW^t(M,L_0,L_1)=\underline{\lim }_H
HW^t(H,L_0,L_1,J)
$$
\end{definition}

(Where we have taken direct limit.)

Taking direct limit of the homology,
we make sure the homology theory is
independent of the
choice of objects (I think, complex structure
and Hamiltonian)
we used to construct it.


\begin{proposition}
\label{proposition-persistence-module}
$t \to HW^t(M,L_0,L_1)$ is a peristence
module $B(M,L_0,L_1).$
\end{proposition}

Finally we can define {\it barcode entropy}.
Fix $\varepsilon>0$,
$t \geq 0$,
$$
b_\varepsilon(M,L_0,L_1,t)=
\# \{\text{ of bars in }B(M,L_0,L_1)
\text{with length $\geq \varepsilon$
and start before $t$}\}
$$
Then
$$
\bar{h}^{HW}(M_0,L_0,L_1)=
\lim_{\varepsilon \to 0} \text{lim sup}_{t
\to \infty}
\frac{\text{log}^+(b_\varepsilon(M,L_0,L_1,t)
}{t}
$$

\medskip\noindent
Consider a contact manifold
$(\Sigma,\lambda,L_0,L_1)$
and $A$ the Lagrangians. (Example raising
question of filling.)

\begin{theorem}[M '24]
\label{theorem-entro-indep-filli}
$\bar{h}^{HW}$ is independent of the filling.
\end{theorem}

\begin{theorem}[M '24]
\label{theorem-bound}
$\bar{h}^{HW}(M,L_0,L_1) \leq
h_{\text{top}}(\alpha)$.
\end{theorem}

\begin{theorem}[M '25]
\label{theorem-restriction}
Consider $(M,L_0,L_1)$.
Let $K$ be a compact topologically transitive
hyperbolic set
for the Reeb flow $\alpha$. Assume $W_\delta
^s (q) \subset \partial L_0$,
$W_\delta^s(p) \subset \partial L_1$. Then
$$
\bar{h}^{HW}(M,L_0,L_1)\geq
h_{\text{top}}(\alpha |_{K})>0.
$$
\end{theorem}
Which says it captures dynamics beyond
unconditional phenomena.
In lower dimensions these tend to coincide,
but in higher dimension we don't
know. This is related to
 ``the sup over hyperbolic sets […]''

Here's a conjecture:
$$
\text{sup}_{L_0,L_1}\bar{h}^{HW}(M,L_0,L_1)=
h_{\text{top}}(\alpha).
$$
\medskip\noindent
{\bf Extra comments.}
One of the aims is to describe topological
entropy $h_\text{top}$
using Floer theory. Theorems by
Çineli-Ginzburg-Gürel show bounds
of topological entropy and barcode entropy
(one of which is for
hyperbolic sets).

There is a notion of  {\it admissible
Hamiltonian and Reeb vector fields}
which is related to some asymptotical
behaviour  ``linear at infinity''.
I understand that admissible vector fields
give
the interesting chords for the Floer homology
construction.

\section{Revisiting cotangent bundles}
\label{section-revisiting-cotangent-bundles}
%% speaker: Mieugel Cueca, KU Leuven
%% place: Symplectic Geometry Joint Seminar,
%% place: IMPA
%% date: September 5, 2025

\noindent
Mieugel Cueca, KU Leuven.
Symplectic Geometry Joint Seminar, IMPA.
September 5, 2025.

\medskip\noindent
{\bf Abstract.} Cotangent bundles provide key
examples of symplectic
manifolds. On the other hand, one can think
of Lie groupoids as generalizations
of manifolds. In this context, Alan Weinstein
constructed their cotangent
bundles and proved that they are so-called
symplectic groupoids. In this talk, I
will recall this construction and explain
what happens when one replaces a Lie
groupoid with a Lie 2 (or n)-groupoid. If
time permits, I will exhibit some of
their main applications. This is joint work
with Stefano Ronchi.

\medskip\noindent

Recall the basic properties of the cotangent
bundle $T^*M$ for a symplectic manifold:
\begin{enumerate}
\item It's a vector bundle.
\item $\left\langle{}\cdot,\cdot\right\rangle{}:TM \otimes
T^*M \to \mathbb{R}_M$
is the dual pairing.
\item $\omega_{\text{can}}\in \Omega^2(T^*M)$
is symplectic.
\item $\mathcal{L}_\varepsilon
\omega_{\text{can}}=\omega_{\text{can}}$,
$\varepsilon$ Euler vector field.
\end{enumerate}

\medskip\noindent
{\bf Main goal.} Reproduce the above for Lie
$n$-groupoids.

For $n=0$ we get the above situation. For
$n=1$ [Duzud-Weinstein], [Prezun],
for $n\geq 2$ ? and we care about $n=2$.

\begin{definition}
\label{definition-groupoid}
A {\it Lie $n$-groupoid}
$\mathcal{G}:\Delta^{\text{op}}\to
\mathsf{Man}$
such that
$$
P_{e,j}:\mathcal{G}_{\ell}\to \Lambda_j^\ell
\mathcal{G}
$$
are surjective submersions $\forall  \ell,j$
are differomorphisms for $\ell > n$.
\end{definition}

\begin{remark}
\label{remark-generators}
This sort of manifolds-valued presheaf
category is generated by
$$
\begin{aligned}
d_i^\ell:\mathcal{G}_\ell \to
\mathcal{G}_{\ell-1}&\text{face maps}
0 \leq  i,j \leq  \ell\\
s_j^\ell:\mathcal{G}_\ell \to
\mathcal{G}_{\ell+1}&  \qquad
\text{degeneracies}
\end{aligned}
$$
\end{remark}

\medskip\noindent
The tangent is a functor, it satisfies
$$
T_\bullet(\mathcal{G})=T_k(\mathcal{G})=
T\mathcal{G}_k
$$
(It looks like $T$ preserves diagrams.)

\medskip\noindent
{\bf Dold-Kan.} The category $\mathsf{SVect}$
of simplicial vector spaces
has objects
$$
\xymatrix{
\mathbb{V}_\bullet&
\mathbb{V}_n\ar[r]&\cdots
\ar[r]&\mathbb{V}_2\ar[r]^{\text{3 arrows}}&
\mathbb{V}_1\ar[r]^{\text{2
arrows}}&\mathbb{V}_0
}
$$
where the $\mathbb{V}_i$ are vector spaces.

There is a functor
$$
\mathsf{SVect}\overset{N}{\to}\{\text{chain
complexes}\geq 0\}
$$
$$
\mathbb{V}_{\bullet}\to N\mathbb{V}=
\left(N_\ell \mathbb{V}\Ker
P_{\ell,\ell},\partial=d_\ell\right)
$$

\begin{theorem}[Dold-Kan]
\label{theorem-Dold-Kan}
That's an equivalence of categories. [Confirm
this!]
\end{theorem}

Those categories are monodial:
$$
(\mathsf{SVect},\otimes),\qquad
(\mathbb{V}_\bullet \otimes \mathbb{W})_\ell
=\mathbb{V}_\ell \otimes \mathbb{W}_\ell
$$
$$
(\mathsf{ch}_{\geq 0},\otimes),\qquad (V
\otimes W)_i=
\bigoplus_{\ell +k=1}V_\ell \otimes W_k
$$
And $N$ is Lax monoidal with Lax structure
given by the
Eilenberg-Zilber map, though we won't explain
the details of this.

There are duals given by internal Hom:
$$
\begin{aligned}
\mathbb{V}^{n*}&=\underline{\Hom}(\mathbb{V},
B^n\mathbb{R})
\end{aligned}
$$
Where the internal Hom is given by
$$
\underline{\Hom}(\mathbb{V},B^n\mathbb{R})
_\ell
=\Hom_{\mathsf{SVect}}(\mathbb{V}
\otimes\Delta_n[\ell],B^n\mathbb{R})
$$
for an object
$\Delta[\ell]=\mathbb{R}[\Delta[\ell]]$.


\medskip\noindent
{\bf Properties.}
\begin{enumerate}
\item $\mathbb{V}^{n*}$ isa simplicial vector
bundle.
\item $N(V^{n*})$ and $N(\mathbb{V})^*[n]$ is
a quasi isomorphism.
\item
$\left\langle{}\cdot,\cdot\right\rangle{}:\mathbb{V}\otimes
\mathbb{V}^{n*}\to
B^n\mathbb{R}$ is non-degenerate on homology.
\item $\mathbb{V}\hookrightarrow
(\mathbb{V}^{n*})^{n*}$
Mont. a eq.
\end{enumerate}

\medskip\noindent
{\bf The vector bundle case.}
$$
\text{Maps}(\Delta[i],\mathcal{G})_k=
\Hom_{\mathsf{SSet}}(\Delta[i]\times\Delta[k]
,\mathcal{G})
$$
\begin{proposition}
\label{proposition-on-an-n-Lie-groupoid}
Let $\mathcal{G}$ be a Lie $n$-groupoid.
\begin{enumerate}
\item $\text{Maps}(\Delta[i],\mathcal{G})$
Lie $n$-groupoids ME $\mathcal{G}$.
\item
$\text{Maps}(\Delta[i],\mathcal{G})_0=
\mathcal{G}_i$.
\item $\text{ev}:\Delta[i]\times
\text{Maps}(\Delta[i],\mathcal{G})\to
\mathcal{G}$.
\end{enumerate}
\end{proposition}

[Staircase looking diagram.]

\begin{definition}
\label{definition-what-is-this}
$\mathcal{G}_\bullet$.
$$
T_i^{n*}\mathcal{G}=\Hom_{\mathsf{SVect}}(1^*
\Pi_{\Delta[i]}(T\mathcal{G}),
B^n\mathbb{R}_{\mathcal{G}_i})
$$
$$
(d_j,F)_{K|d_j\mathcal{G}}(x^a)=(F_k)
|_{\mathcal{G}}(x^{\delta_ja}).
$$
\end{definition}

\begin{proposition}
\label{proposition-properties-of-that}
$\mathcal{G}$ Lie $n$-groupoid, then
$T^{n*}\mathcal{G}$ satisfy
\begin{enumerate}
\item is a vector bundle $n$-groupoid
\item dual to $T\mathcal{G}$
$$
\left\langle{}\cdot,\cdot\right\rangle{}:T\mathcal{G} \otimes
T^{n*}\mathcal{G}
\to B^n\mathbb{R}_{\mathcal{G}}
$$
non-degenerate on homology.
\item $n$-shifted symplectic
$$
T^{n*}_n\mathcal{G}\overset{p}{\to}
T^*\mathcal{G}_n
$$
and $p^*\omega_{\text{can}}$.
\end{enumerate}
\end{proposition}

[More computations I missed]

\section{Holomorphic extensions of s-proper
Lie groupoids}
\label{section-holom-exten-s-prope-lie}
%% speaker: Rui L. Fernandes, Instituto Superior
%% speaker: Técnico (Lisboa) / University of
%% speaker: Illinois at Urbana-Champaign
%% place: Symplectic Geometry Seminar, IMPA
%% date: September 17, 2025

\noindent
Rui L. Fernandes, Instituto Superior Técnico
(Lisboa)/University of Illinois
at Urbana-Champaign.
Symplectic Geometry Seminar, IMPA.
September 17, 2025.

\medskip\noindent
{\bf Abstract.} Every smooth manifold admits
a compatible analytic
structure, and a classical result of
Whitney–Bruhat states that any analytic
manifold has a holomorphic extension. Lie
groups also admit compatible analytic
structures, and another classical result, due
to C. Chevalley, shows that any
compact Lie group has a holomorphic extension
to a complex Lie group. D.
Martínez Torres has shown that any proper Lie
groupoid admits a compatible
analytic structure. I will discuss an
extension of the classical results of
Whitney–Bruhat and Chevalley, establishing
that any s-proper Lie groupoid has a
holomorphic extension. This talk is based on
recent joint work with Ning Jiang
(arXiv:2508.18036).

\medskip\noindent

\begin{theorem}[Whitney-Beuhat]
\label{theorem-Whitney-Beuhat}
If $M$ is analytic there exists a complex
manifold  $M_\mathbb{C}$
together with an analytic map $i: M \to
M_\mathbb{C}$
which is totally real ($T_HM_{\mathbb{C}}=TM
\oplus J(TM)$)
and
\begin{enumerate}
\item For every complex manifold $X$ and
every
analytic map $\phi:N \to X$ there exists
$\phi^*:U \to X$
holomorphic contained open $M \subset U
\subset M_\mathbb{C}$.
\item If $\psi:V \to X$ holomorphic on $M
\subset V \subset M_\mathbb{C}$ then
$\psi=\phi^*$ where $\phi=\psi \circ i$ on a
possibly
smaller open contained in $U \cap V$.
\end{enumerate}
\end{theorem}

\begin{theorem}[Chevalley]
\label{theorem-Chevalley}
If $G$ is a Lie group, there exists a complex
Lie group $G_{\mathbb{C}}$ and a morphism
$i:G \to G_\mathbb{C}$
satisfying the following universal property.
For every complex Lie group $H$ and morphism
$\phi:G \to H$ there exists a unique
holomorphic map $\phi^*:G_\mathbb{C} \to H$
such that
$\phi=\phi^* \circ i$.
\end{theorem}

Here's the construction of $G_\mathbb{C}$:

$$
\xymatrix{
N\subset \tilde{G}\ar[r]\ar[d]&G^*\ar[d]\\
\tilde{G}/N \simeq
G\ar[r]_{i}&G_\mathbb{C}=G^* /\overline{i^*N}
}
$$
where $\tilde{G}$ is the universal cover
group of $G$,
 $G^*$ is the group that integrates to
the complexification of the Lie algebra of
$G$,
i.e.
$\text{Lie}(G^*)=\mathfrak{g}_\mathbb{C}$
and $\overline{i^*(N)}$ is the smallest
closed
normal complex Lie subgroup of $G^*$
containing $i^*(N)$.

\begin{example}
\label{example-complexification-groupoid}
$$
\xymatrix{
\widehat{\text{SL}_2(\mathbb{R})}\ar[d]
\ar[dr]\\
\text{SL}_2(\mathbb{R})\ar@{^{(}->}[r]&
\text{SL}_2(\mathbb{C})
}
$$
which is a simple case,
but consider instead
$\tilde{\text{SL}_2}(\mathbb{R})\times
\mathbb{R}$
and the normal subgroup
$N=\left\langle{}a\right\rangle{}\times \left\langle{}\lambda\right\rangle{}$
for irrational $\lambda$.
Then we obtain
$$
\xymatrix{
\widehat{\text{SL}_2(\mathbb{R})}
\times\mathbb{R}\ar[r]\ar[d]
&\text{SL}_2(\mathbb{C})\times\mathbb{C}
\ar[d]\\
\widehat{\text{SL}_2(\mathbb{R})}
\times\mathbb{R}/N\simeq
G\ar[r]&
G_\mathbb{C}
}
$$
and $G_\mathbb{C}$ is 3 complex dimensions!
It is not
$\text{SL}_2(\mathbb{C})\times\mathbb{C}$.
\end{example}

\section{Lagrangian blow-up and blow-down for
4-dimensional symplectic manifolds
}
\label{section-blow-up-down}
%% speaker: Misha Verbitsky, IMPA
%% place: Geometric Structures Seminar, IMPA
%% date: October 23, 2025

\noindent
Misha Verbitsky, IMPA.
Geoemtric Structures Seminar, IMPA.
October 23, 2025.

\medskip\noindent
{\bf Abstract.} The usual (complex geometric)
blow-up has a symplectic
version, called the symplectic cut. I will
introduce a variant of this
construction, which is valid only in
symplectic category, called ``Lagrangian
blow-up'', and its inverse, called Lagrangian
blow-down. Lagrangian blow-down
takes a Lagrangian sphere in a 4-dimensional
symplectic manifold and contracts
it to a symplectic orbifold with a double
point. Given a symplectic orbifold
with a double point, Lagrangian blow-up
produces a symplectic manifold, and this
construction is inverse to the Lagrangian
blow-down. I will explain why these
constructions are functorial, that is,
defined on the corresponding symplectic
Teichmuller spaces. This is used to prove an
orbifold counterexample to a famous
conjecture, sometimes attributed to
Donaldson, who asked whether any symplectic
form on a K3 surface is compatible with a
Kahler structure. I will prove that
this is false for orbifolds: some K3
orbifolds admit symplectic forms not
compatible with any Kahler structure. The
same argument is used to produce a
countable family of Lagrangian spheres in a
K3 surface which are not Lagrangian
isotopic to special Lagrangian spheres. This
is a joint work with Michael Entov.



\medskip\noindent


Consider the total space of the cotangent
bundle of $\mathbb{C}P^1$.
It's $\mathcal{O}(-2)$.
Let $\gamma \in H^{0}(\mathcal{O}(2))$.
$\gamma$ defines a fiberwise linear function
on $T^*\mathbb{C}P$.

The blow-up of $\mathbb{C}^2/\pm 1$, which is
the orbifold $\mathbb{C}^2$ with
 a doble point, is the cotangent bundle $T^*
 \mathbb{C}P^1$.
That is,
$T^*\mathbb{C}P^1 \to \mathbb{C}^2/\pm 1$.
This is called {\it crepant resolution}.
Consider the $dx \wedge dy$, which is
holomorphically symplectic
on $\mathbb{C}^2$ and descends to the
quotient $\mathbb{C}^2/\pm 1$
as $\Omega_0$.

\begin{proposition}
\label{proposition-}
$\pi:T^*\mathbb{C}P^1 \to \mathbb{C}^2/\pm
1$.
Then $\pi^*\Omega_0$ is holomorphically
symplectic.
\end{proposition}

\begin{proof}
Actually, it gives the standard symplectic
form on $T^*\mathbb{C}^2$.
\end{proof}

\medskip\noindent
Now we shall define the Langrangian blow-up.
It's a procedure that starts with a
symplectic manifold
with double points and outputs a symplectic
manifold
without double points.

\begin{definition}[Lagrangian blow up]
\label{definition-lagrangian-blow-up}
Let $(B,\omega)$ be the standard ball in
$\mathbb{R}^n$,
$dz_1 \wedge dz_2 +dz_3 \wedge dz_4$.
Consider the symplectic orbifold
$(B,\omega)/\pm 1$.
Consider $M$ be a symplectic orbifold
with double point singularities.
Darboux coordinates in a neighbourhood of $0$
with $(B,\omega)/\pm 1$ resolve
singularities.
\end{definition}

This definition has the caveat that there ar
e lots of choices involved.
We shall later that these choices do not
alter the result,
namely Theorem
\ref{theorem-lagrangian-blow-up-well-defined}
.

To define the Lagrangian blow-down we first
recall

\begin{theorem}[Weinstein's Lagrangian
neighbourhood]
\label{theorem-weins-lagra-neigh}
Let $X \subset M$ be a compact Lagrangian
submanifold in $(M,\omega)$.
Then there exists a neighbourhood $U$ of $X
\subset M$
which is symplectomorphic to a neighbourhood
of $X$
in $X \subset T^*M$.
\end{theorem}

\noindent
Let $S \subset M$ be a Lagrangian sphere
(diffeomorphic to sphere
and Lagrangian).

\begin{definition}
\label{definition-lagrangian-blow-down}
The {\it Lagrangian blow-down} is obtained
from $M$
by removing $W \subset M$ and gluing $W_0$ in
its place,
where $W_0$ is such that

$$
\xymatrix{
W
\ar@{^{(}->}[r]\ar[d]&T^*\mathbb{C}P^1\ar[d]
\\
W_0\ar@{^{(}->}[r]&\mathbb{C}^2/\pm 1.
}
$$
\end{definition}

\medskip\noindent
Now we recall Teichmüller spaces.

Let $\Gamma(\Lambda^2M)$ be the space of all
2-forms on a manifold $M$.
The space of symplectic forms is a subspace
of the closed 2-forms on $M$,
which in turn is a subspace of
$\Gamma(\Lambda^2M)$.
Equip this space with the $C^\infty$
topology.
The {\it Teichmüller space} is
$\text{Symp}/\text{Diff}_0$.

Two symplectic structures are called {\it
isotopic} if they lie
in the same orbit of $\text{Diff}_0$.
The {\it period map}
$\text{Per}:\text{Teich}_s \to
H^{2}(M,\mathbb{R})$
maps a symplectic structure to its cohomology
class.

\begin{theorem}[Moser, 1965]
\label{theorem-period-map-is-local-diffeo}
The period map is a local diffeomorphism.
\end{theorem}

\noindent
This implies that $\text{Teich}_s$ is smooth.

\begin{theorem}[Moser's trick]
\label{theorem-mosers-trick}
Let $\omega_t$, $t \in S$ be a smooth family
of symplectic
structures, parametrized by a connected
manifold $S$.
Assume that the cohomology class
$[\omega_t]$ is constant.
Then all $\omega_t$ are diffeomorphic.
\end{theorem}

\noindent
This means that the fibers of $\text{Symp}\to
H^{2}(M)$
are the orbits (of the action of
$\text{Diff}_0$, I suppose).

\medskip\noindent
Consider a Lagrangian submanifold $L \subset
(M,\omega)$.
Let us assume for simplicity (though the
results
can also be obtained without this) that
$b_1(L)=0$,
i.e. first cohomology group is zero.

Let $\text{Symp}_L$ be the set of symplectic
forms vanishing on $L$,
and $\text{Diff}_0^L$ the connected component
of diffeomorphisms
preserving $L$.
Define
$\text{Teich}_L=\text{Symp}_L/\text{Diff}
_0^L$.

\begin{theorem}
\label{theorem-}
Let $V \subset H^{2}(M,\mathbb{R})$ be the
space of all
cohomology classes on $M$ which vanish on
$L$,
and $\text{Per}_L:
\text{Symp}_L/\text{Diff}^L_0 \to V$
take $(M,\omega)$ to the cohomology class of
$\omega$.
Then Per is a local diffeomorphism.
\end{theorem}

\begin{proof}
We used a version of Moser's trick: for
$\omega_t$ a family
of symplectic forms with $\omega_t|_{L}=0$,
with $[\omega_t]$ constant, then $\omega_t$
are $\text{Diff}_0^L$-isotropic.
\end{proof}

\noindent
As a corollary we obtain that this
Teichmüller space
classifies Lagrangian submanifolds:

\begin{corollary}
\label{lemma-techm-class-lagra-subma}
Assuma that $L$ and $L'$ are Lagrangian
submanifolds in $(M,\omega)$,
and $\varphi \in \text{Diff}_0(M)$ a smooth
isotopy such that
$\varphi(L)=L'$. Then $(\omega,L)$ and
$(\varphi^*\omega,L)$
represent the same point in $\text{Teich}_L$
if and only if
$L$ and $L'$ are Lagrangian isotopic in
$(M,\omega)$.
\end{corollary}

\noindent
Again: points of this Teichmüller space are
isotopy classes
of Lagrangian submanifolds.

\medskip\noindent
\begin{theorem}
\label{theorem-blow-down-diffeomorphism}
Let $L$ be a Lagrangian 2-sphere in a K3
surface  $M$
and $\hat{M}$ the Lagrangian blow-down of
$M$.
The blow-up/blow-down
is a diffeomorphism $\text{Teich}_L(M) \to
\text{Teich}(\hat{M})$.
\end{theorem}

\medskip\noindent
Now we prove the theorem which asserts that
the Lagrangian
blow-up is independent of choices.

\begin{theorem}
\label{theorem-lagrangian-blow-up-well-defined}
Let $(\hat{M},\hat{\omega})$ be a orbifold.
Then its blow-up
$(M,L)$ defines a point in $\text{Teich}_L$
independent of the choices.
\end{theorem}

\begin{proof}
By taking local Darboux coordinates, and then
Moser's trick.
\end{proof}

\noindent
Actually, the same happens for blow-downs

\begin{theorem}
\label{theorem-same-for-blow-down}
Let $(M,L,\omega)$ be a 4-manifold with a
Lagrangian sphere $L$, and
$(\hat{M},\hat{\omega)}$ be its Lagrangian
blow-down.

Then the corresponding point in
$\text{Teich}_{\hat{M}}$ is
independent from the choices made.
\end{theorem}

\begin{proof}
Similar.
\end{proof}

\medskip\noindent
A conjecture by Donaldson ways that all
symplectic structures on a K3
surface are compatible with a Kähler
structure. This conjecture
(although it's commonly believed to be true)
is still open.
When trying to prove that this would hold for
orbifolds,
Misha and collaborators realised it's
actually false.

Note: Seiberg-Witten invariants may be used
to tell whether a symplectic
form is compatible with a complex structure.

\begin{theorem}
\label{theorem-conje-false-orbif}
There exists an orbifold symplectic K3
surface $M$ with a single
double point not admitting an orbifold Kähler
structures.
\end{theorem}

\begin{proof}
\noindent
Step 1. Consider
$$
\text{Teich}_{\text{Kähler-symplectic}}
(\hat{M})
=\{\eta \in H^{2}(\hat{M}):\eta^2>0\}.
$$
As follows from [Amerik-V., 2005], the
Teichmüller space of symplectic structures
compatible
with a hyperkähler structure is Hausdorff and
connected
(the same argument works for Kähler K3
orbifolds).

Essentially, if there wasn't any orbifold
as required, the Teichmüller space
classifying
Lagrangian submanifolds would be connected,
which
is impossible by [Seidel 2000].
\end{proof}

\noindent
In fact, Seidel's result is essentially
mirror symmetry.

\medskip\noindent
Now we discuss special Lagrangian
submanifolds.

\begin{definition}
\label{definition-phase}
Let $(M,I,\omega,\Omega)$ be a Calabi-Yau
manifold, where
$\Omega \in \Lambda^{n,0}_I(M)$ is
nondegenerate and $\omega$ is Kähler.
Let $L$ be a Lagrangian submanifold.
Define the {\it phase} as
$$
\begin{aligned}
\text{phase}: L &\longrightarrow \Lambda^nL
\otimes \mathbb{C} \\
x &\longmapsto \Omega|_{T_xL}.
\end{aligned}
$$
\end{definition}

\noindent
A better definition is:
$$
\begin{aligned}
\text{phase}: L &\longrightarrow
S^1=\text{U}(1) \\
x &\longmapsto \frac{\Omega|_{T_xL}}{|\Omega
T_xL|}.
\end{aligned}
$$

\begin{definition}
\label{definition-special-lagrangian}
A special Lagrangian submanifold is the phase
is constant.
\end{definition}

\noindent
Special Lagrangian submanifolds have
deformation space of
dimension 1, it's unobstructed. They are
calibrated, so minimal.

\begin{theorem}
\label{theorem-lagrangians-of-k3}
$L \subset$ K3 is Lagrangian is isotopic to a
special Lagrangian
if and only if its blow-down is Kähler type.
\end{theorem}

\medskip\noindent
One of the basic results about special
Lagrangian
submanifolds is

\begin{lemma}[Hitchin]
\label{lemma-hitchin}
If $(M,\Omega)$ is holomorphically symplectic
and $X \subset M$
is Lagrangian with respect to
$\text{Re}\Omega$ and $\text{Im}\Omega$
then $X$ is complex.
\end{lemma}

\noindent
As corollaries:

\begin{corollary}
\label{lemma-corollary-hitchin}
Let $(M,I,J,K,g)$ be a hyperkähler manifold
of real dimension $4n$,
considered as a Kähler manifold
$(M,I,\omega_I)$ and $\phi:\Omega^n \in
\Lambda^{2n,0}(M,I)$ the corresponding
holomorphic volume form.
Consider a Lagrangian submaniold $S \subset
(M,\omega_I)$
such that $a \omega_J + b\omega_K|_S=0$ for
some nonzero real numbers
$a,b \in \mathbb{R}$ such that $a^2+b^2=1$.
Then $S$ is special Lagrangian, and,
moreover,
it is holomorphic with respect to the complex
structure
$-a K+b J$.
\end{corollary}

\begin{corollary}
\label{lemma-corollary2-hitchin}
Let $\omega_I,\omega_J$ and $\omega_K$ on
$M$, with $\dim M=4$.
Let $S \subset (M,\omega_I)$ be Lagrangian.
Then $S$ is special Lagrangian if and only if
$S$ is complex analytic
on $a J + b K$ for two nonzero numbers  $a,b
\in \mathbb{R}$ such that
$a^2+b^2=1$.
\end{corollary}

\noindent
Now an auxiliary claim:

\begin{lemma}
\label{lemma-k3-type}
$\Omega$ is holomorphic symplectic
\end{lemma}

\begin{theorem}
\label{theorem-duplicate-2}
Let $S\subset M$ be a Lagrangian sphere in a
K3 surface.
Then  $S$ is isotropic to a special
Lagrangian if and only if
$M$ is of Kähler type.
\end{theorem}

\section{From higher Lie groupoids to higher
Lie algebroids}
\label{section-highe-group-highe-algeb}
%% speaker: Matias del Hoyo, UFRJ
%% place: Symplectic Geometry Seminar, IMPA
%% date: October 29, 2025

\noindent
Matias del Hoyo, UFRJ.
Symplectic Geometry Seminar, IMPA.
October 29, 2025.

\medskip\noindent
{\bf Abstract.}
This is a report of an ongoing collaboration
with A. Cabrera, where we develop a
higher Lie theory. Starting with simple,
concrete examples, we present the
differentiation functor in a geometric,
explicit way, using simplicial methods
and differential graded algebras. We describe
how our construction extends the
classical Lie functors, relate our project
with previous works, and show a van
Est theorem linking the cohomology at the
global and infinitesimal levels.

\medskip\noindent

Upshot: Integration-differentiation relation
between Lie groups and Lie algebras
in higher versions. Groupoids are the
counterparts of Poisson manifolds, and in
the next level there are 2-groupoids and
Courant algebroids. How to integrate a
2-groupoid to a Courant algebroid? This is an
open question which motivated this
work. See arXiv
\href{http://arxiv.org/abs/2309.14105}%
{2309.14105}.

\medskip\noindent
Let $M$ be a smooth manifold.

\begin{remark}
\label{remark-smoot-funct-ring-fully-faith}
$M \mapsto C^\infty M$ is fully faithful.
\end{remark}

\noindent
Which allows us to think of manifolds as
algebraic objects.
Notice that
$\Omega^kM=\{k\text{-forms}\}
=\Lambda^k_{C^\infty M}\text{Der}(C^\infty
M,C^\infty M)$,
which is (algebraically) locally spanned
by the exact 1-forms.

Recall that the exterior differential is
the only operator $d:\Omega^kM \to
\Omega^{k+1}M$
such that $d(\alpha \wedge \beta)=d\alpha
\wedge \beta +
(-1)^{|\alpha|}\alpha \wedge d\beta$
and that $df(X)=Xf$.

This makes $(\Omega M,\wedge,d)$ a
differential graded algebra.
Recall that $\text{dom}(X,Y)=X\omega(Y)-Y
\omega(X)-\omega([X,Y])$.
We also have de Rham cohomology.

Consider a submanifold $S \subset M$. Recall
that $\mathcal{I}_S/\mathcal{I}^2_S$ is
essentially
the sections of the normal bundle.
In particular, for $M \subset M \times M$
we can see that the embedding vanishes on the
diagonal,
so there is no linear part in the Taylor
expansion,
and the second term is in fact the de Rham
differential
$df=\pi_1^*f-\pi_2^*f \text{ mod
}\mathcal{I}^2_{\Delta(M)}$.

In fact, this construction can be extended to
a simplicial manifold
(which I think is just a presheaf valued on
category of manifolds)
using projections and diagonals,
$$
\xymatrix{
\cdots & M\times M \times M
\ar @< 3pt> [r]
\ar @< 0pt> [r]
\ar @<-3pt> [r] & M \times M \ar @< 2pt> [r]
\ar @<-2pt> [r]& M.
}
$$
To this we can apply the $C^\infty$ functor
and obtain
a cosimplicial ring (not a dga?).
The idea here
is that de Rham differential is an
alternating sum of pullbacks;
as other differentials will also be.

\medskip\noindent
Now observe that for a finite-dimensional
vector space $V$,
$$
\left\{\substack{\text{Lie bracket} \\
V \times V
\xrightarrow{[\cdot,\cdot]}V}\right\}
\leftrightarrow
\left\{\substack{\text{differential on} \\
\Lambda V^*}\right\},
$$
making $\Lambda V^*$ into a differential
graded algebra.

Recall that given a Lie algebra
$\mathfrak{g}$,
we have $(\Lambda^\bullet \mathfrak{g}^*
,d)=\text{CE}(\mathfrak{g})$,
the Chevalley Eilenberg algebra,
which gives Lie algebra cohomology.

\medskip\noindent
[Missing part on Simplicial approach,
exponential map,
Simplicial manifold nerve]

\medskip\noindent
Lie algebroid: $\Delta \to M$ vector bundle
with $[\cdot,\cdot]$ on $\Gamma A$,
$A \xrightarrow{\rho}TM$ Leibniz.

\begin{remark}
\label{remark-lie-algebroids-differentials}
$$
\left\{\substack{\text{Lie algebroids on} \\
E \to M}\right\}
\leftrightarrow
\left\{\substack{\text{differentials} \\
\text{on }\Lambda\Gamma E^*}\right\}.
$$
\end{remark}

\noindent
This gives a differential graded algebra
$$
\xymatrix{
\Delta\ar[r]&M\ar[r]&\text{CE}(A)
}
$$
and
$$
\xymatrix{
C^\infty M\ar[r]^{\tilde{\rho}}&\Gamma
E^*\ar[r]^{\widetilde{[\cdot,\cdot]}}
&\Gamma E^* \wedge \Gamma E^*\ar[r]&\cdots.
}
$$

From Lie groupoids to Lie algebroids:
Lie functor/differentiation.

$$
\Omega(G)=\Gamma \Lambda^\bullet
\underbrace{\supset}
_{\text{sub dga}}
\Gamma \Lambda^\bullet T^s
\xrightarrow{-/G}\text{CE}(A_G).
$$
There's the following ecosystem around Lie
groupoids:
$$
\xymatrix{
&\text{Lie groupoids}
\ar[dl]\ar@{->>}[d]\ar[dr]^{\text{nerve}}\\
\text{Lie algebroids}&
\substack{\text{Differentiable stacks} \\
\supset \text{orbifolds}}
&\text{simplicial manifolds.}
}
$$

\medskip\noindent
Lie groupoids are models for stacks,
which are geometric objects solving
classification
problems with objects identified up to
isomorphism.
Relaxing the identification to a weaker
notion
of equivalence we must introduce higher
versions of stacks.

Moving higher,
$$
\xymatrix{
\text{Lie groupoid}\ar[r]\ar[d]_{\text{Lie}}
&\text{Simplicial manifolds}\ar[d]^?\\
\text{Lie
algebroids}\ar[r]&\text{dg-manifolds}
}
$$
dg algebra is a  dga $(A=\bigoplus_{k=1}^n
A_k \to M,\wedge,d)$
starting at $C^\infty M$, and continuing
until you get to
the sections, namely
$$
\xymatrix{
C^\infty M \ar[r]_d & A_2 \ar[r]& \cdots &
\cong \Gamma A.
}
$$
[dH-Cabrera] provide the solution to this
problem,
namely what to put in the right vertical
arrow:

\begin{theorem}[dH-Cabrera]
\label{theorem-dh-cabrera}
$$
\frac{(C^\infty_N(G),\cup,\sum(-1)A^*_i)}{J+
\delta
J}
=\text{CE}(A_G)
$$
where $A_G$ is a higher Lie algebroid
and $J$ is defined by: $f$ is in $J$ if it
has higher vanishing.
\end{theorem}


Groupoids are the counterparts of Poisson
manifolds,
and in the next level there are 2-groupoids
and Courant algebroids.
How to integrate a 2-groupoid to a Courant
algebroids?
This is an open question which motivated this
work.

\section{Rafael Fernandes, Differential Geometry
Seminar}
\label{section-fernandes-today}
%% speaker: Rafael Fernandes, UC Santa Cruz
%% place: Differential Geometry Seminar, IMPA
%% date: September 11, 2026
\medskip\noindent
{\bf Speaker.}

Rafael Fernandes, UC Santa Cruz.

\medskip\noindent
{\bf Date.}

September 11, 2026.

\medskip\noindent
{\bf Seminar.}

Differential Geometry Seminar, IMPA.

\medskip\noindent {\bf Title.} On the growth
rate of Reeb orbits on star-shaped
hypersurfaces


\medskip\noindent
{\bf Abstract.}

In this talk, I will discuss the growth rate
of Reeb orbits with respect to their periods
on certain contact manifolds. We focus on
fiberwise star-shaped hypersurfaces in the
cotangent bundle of a closed manifold. This
class of dynamical systems includes, for
example, geodesic flows. I will explain how
the string topology of the base manifold
implies that the number of Reeb orbits with
period at most grows at least like the prime
numbers, that is, on the order of . The
topological condition is the existence of a
certain non-nilpotent homology class of the
free loop space (with respect to the
Chas–Sullivan product). This is joint work
with João Pering.

\medskip\noindent
The following theorem links topology and
dynamics:

\begin{theorem}[Bangart-Hingston]
\label{theorem-bangart-hingston}
Let $(M^n,g)$, $n \geq 2$, closed.
If $\pi_1(M)$ is infinite and abelian,
then
$$
\lim_{t \to \infty} N(t)
\frac{\operatorname{log}(t)}{t}>0
\iff
N(t) \geq C \frac{t}{\operatorname{log} t}.
$$
\end{theorem}

\noindent
{\bf Idea.} Lustenik-Shcnrelmann
theory + construction of
some higher homotopy groups of free loop
spaces.

$\phi_s:S^\bullet M \to S^\bullet M$, bijective
correspondence between geodesics of length
up to $t$ and closed orbits of $\phi_s$
with period up to $t$.

\medskip\noindent
Note. $S^\bullet M \subset T^\bullet M$.

\medskip\noindent
Question: does the growth hold for more
general higher surfaces of contact type of
$T^\bullet M$?

\begin{definition}
\label{definition-fiberwise-star-shaped}
$D \subset T^*M$ is {\it fiberwise
star-shaped} if $D \cap T^*_+M$ is
starshaped and $\partial D = \Sigma$ is
transverse to the radial direction.
\end{definition}

\noindent
$$
\Lambda M
=
\coprod_{\eta \in H_1(M,\mathbb{Z})}
\Lambda_\eta M,
$$
$$
H_*(\Lambda M) =
\bigoplus_{\eta \in H_1(M,\mathbb{Z})}
H_*(\Lambda_\eta M).
$$

\begin{definition}
\label{definition-h-inf-non-nilpotent}
$\beta \in H_*(\Lambda_\eta M)$
is called {\it $H_1$-infinity non-nilpotent} 
if $\beta^m \neq  0$ for all $m \geq  1$
and $\eta \neq  \in H_1(M,\mathbb{Z})/Tor$,
i.e. $\eta $ is not torsion.
\end{definition}

\begin{theorem}[F-Pering]
\label{theorem-f-pering}
Let $M^n$ be a closed manifold,
$n>2$. If there exists $\beta$
$H_1$-infinity non-nilpotent, then
($\eta=\eta(\beta)$)
$$
\operatorname{lim inf}_{t \to \infty}
N^\eta_\Sigma(t)
\frac{\operatorname{log}t}{t}>0
$$
for all $\Sigma$ fiberwise star-shaped.
\end{theorem}

\begin{example}
\label{example-f-pering}
\begin{enumerate}
\item $S^1 \times M^n$, $n>1$.
\item $M \times P$, where $M,P$ satisfy the
theorem.
\item $H$-spaces within $\pi_1$ infinite.
\item Principal $S^1$-bundles
$S^1 \to M \to B$ within
$\pi_0(B) = 0$.
\item $M$ that admit a homologically non
trivial $S^1$-action.
\end{enumerate}
\end{example}

\medskip\noindent
{\bf Symplectic homology and string
topology.} $W^{\partial M}$ Liouville domain
with $c_1(W) =0$.

Very roughly, symplectic homology is…

\begin{theorem}[A-S]
\label{theorem-as}
$$
SH_*(T^*M) \simeq
H_*(\Lambda M)
$$
is an algebra isomorphism.
\end{theorem}

\medskip\noindent
{\bf Local symplectic homology.}
Let $X$ be an isolated 1-periodic orbit
of a hamiltonian $H$ in the completion of
$W$, $\hat{W}$.



\section{Closed magnetic geodesics with every
low energy: existence and localization}
\label{section-magnetic-geodesics-low-energy}
%% speaker: Valerio Assenza
%% place: Sala 236
%% date: October 27, 2026, 15:30

\noindent
Valerio Assenza.

Ter 27 out 2026, 15:30.

SALA 236.

\medskip\noindent
{\bf Abstract.} A existência de órbitas
periódicas em um nível de energia
determinado é um problema clássico e bem
conhecido na dinâmica hamiltoniana. Nesta
palestra, vamos nos concentrar nos sistemas
hamiltonianos que
surgem da dinâmica magnética. Suas
trajetórias se projetam em geodésicas
magnéticas, que descrevem o movimento de uma
partícula carregada em uma variedade
riemanniana na presença de um campo
magnético estático.

Nosso principal resultado diz que, em uma
variedade riemanniana fechada, se a
intensidade magnética $b$ não for
identicamente zero, então todo nível de
energia suficientemente pequeno possui uma
geodésica magnética fechada e contraível,
localizada próxima ao conjunto global de
máximos de $b$. Essa localização leva a
vários resultados de existência e
multiplicidade, inclusive em variedades não
compactas. A demonstração usa uma abordagem
variacional que envolve novas técnicas de
interesse por si só. Este é um trabalho em
conjunto com Leonardo Macarini e Gabriele
Benedetti.

\medskip\noindent
{\bf Sources of these notes.} Photographs of
the slides and an automatic transcript of
the audio. Statements follow the slides;
remarks marked as spoken come from the
transcript. Small slips on the slides are
corrected and flagged in remarks.

\medskip\noindent
{\bf Magnetic dynamics.} Let $(M,g)$ be a
Riemannian manifold, with inner product
$\left\langle{} \cdot,\cdot \right\rangle{}$
and Levi-Civita connection $\nabla$. A
\emph{magnetic form} is a closed 2-form
$\beta \in \Omega^2(M)$, so $d\beta = 0$.
Through the metric, $\beta$ determines a
bundle map
$$
B : TM \to TM,
$$
the \emph{Lorentz force}, by
$$
\left\langle{} B(v),w \right\rangle{}
= \beta(v,w),
$$
for all $v,w$ in the same tangent space.
Since $\beta$ is alternating, $B$ is
skew-adjoint at every point.

A non-constant curve
$\gamma : \mathbb{R} \to M$ is a
\emph{magnetic geodesic}, or a
\emph{$(g,\beta)$-geodesic}, if
\begin{equation}
\label{equation-magnetic-newton}
\nabla_{\dot\gamma}\dot\gamma
= B(\dot\gamma).
\end{equation}

Equation (\ref{equation-magnetic-newton}) is
Newton's law for a charged particle under the
Lorentz force. Because $B$ is skew-adjoint,
$$
\frac{d}{dt}
\left\langle{} \dot\gamma,\dot\gamma
\right\rangle{}
= 2 \left\langle{} B(\dot\gamma),
\dot\gamma \right\rangle{}
= 0,
$$
so the speed $|\dot\gamma|$ is constant. In
other words, the kinetic energy
$$
K : TM \to \mathbb{R},
\quad
K(x,v) = \frac{1}{2}|v|^2,
$$
is constant along solutions. We call $m =
|\dot\gamma|$ the \emph{speed} of $\gamma$.

\begin{example}
\label{example-magnetic-r3}
On flat $\mathbb{R}^3$ write
$$
\beta = b_1\, dx\wedge dy
+ b_2\, dy\wedge dz
+ b_3\, dx\wedge dz,
$$
with $b_i \in C^\infty(\mathbb{R}^3)$. The
form $\beta$ corresponds to a vector field
$\vec{B}$ with divergence zero, and in
coordinates
$$
B(v) =
\begin{pmatrix}
0 & -b_3 & b_2 \\
b_3 & 0 & -b_1 \\
-b_2 & b_1 & 0
\end{pmatrix}
\begin{pmatrix}
v_1 \\ v_2 \\ v_3
\end{pmatrix}
= \vec{B} \times v.
$$
Equation (\ref{equation-magnetic-newton})
becomes the classical
$$
\ddot\gamma = \vec{B} \times \dot\gamma.
$$
(The signs and the order of the components
are as on the slide.)
\end{example}

\begin{example}
\label{example-magnetic-surface}
Let $(M,g)$ be an oriented surface, with
area form $\operatorname{vol}(g)$ and
rotation by $\pi/2$ denoted
$J : TM \to TM$, so that
$$
\operatorname{vol}(g)(v,w)
= \left\langle{} Jv,w \right\rangle{}.
$$
Every 2-form is $\beta = b
\operatorname{vol}(g)$ for a unique function
$b \in C^\infty(M)$, the \emph{magnetic
function}. Then
$$
\left\langle{} Bv,w \right\rangle{}
= b \left\langle{} Jv,w \right\rangle{},
$$
so $B = bJ$ and
(\ref{equation-magnetic-newton}) reads
$$
\nabla_{\dot\gamma}\dot\gamma
= b J\dot\gamma.
$$
Geometrically, $\gamma$ is a
$(g,b)$-geodesic with speed $m$ if and only
if $|\dot\gamma| = m$ and its geodesic
curvature is
$$
k^g(\gamma(t)) = \frac{b(\gamma(t))}{m}.
$$
So on surfaces the problem is that of
prescribing geodesic curvature.
\end{example}

\medskip\noindent
{\bf Low and high speed.} The equation is
not homogeneous in $\dot\gamma$:
reparametrizing a solution with speed $m$ to
speed $\lambda m$ does not give a solution.
So the dynamics changes with the speed. For
large speed the curvature $b/m$ is small and
magnetic geodesics look like Riemannian
geodesics. For small speed the curvature is
large and trajectories curl back near their
starting point: the Lorentz force pulls the
particle back. The slides illustrate this
with a picture, and the speaker described
the Mañé critical value below as the
threshold between the two regimes. If
$\beta = 0$ the speed only reparametrizes
the same geodesic, so nothing changes with
$m$.

Let $\widetilde{M} \to M$ be the universal
cover and $\widetilde\beta$ the lift of
$\beta$. The \emph{Mañé critical value} is
\begin{equation}
\label{equation-mane-critical-value}
c(g,\beta) = \frac{1}{2}
\inf_{\theta}
\sup_{\widetilde{M}}
\|\theta\|^2
\in [0,+\infty],
\end{equation}
where the infimum runs over the 1-forms
$\theta$ on $\widetilde{M}$ with
$d\theta = \widetilde\beta$. If
$\widetilde\beta$ is not exact, put
$c(g,\beta) = +\infty$.

From the definition, $c = 0$ exactly when
$\beta = 0$, and $c < +\infty$ exactly when
$\widetilde\beta$ has a bounded primitive.
The slides record two facts about energy
levels $s$:
\begin{enumerate}
\item If $s > c(g,\beta)$ and
$\pi_1(M)$ is abelian, then after lifting to
$\widetilde{M}$ the flow on the level is
conjugate, up to time reparametrization, to
the geodesic flow of a Finsler metric.
\item For $s \le c(g,\beta)$ many basic
questions are open, among them the
existence of closed magnetic geodesics and
of conjugate points.
\end{enumerate}

\begin{remark}
\label{remark-speed-energy-mane}
The slides compare $c$ sometimes with an
energy level $s$ and sometimes with a speed
$m$. With $K = m^2/2$ these differ; the
notes keep the wording of each slide. Which
normalization the papers use was not
checked.
\end{remark}

\begin{example}
\label{example-mane-standard-models}
Three models on closed surfaces, each with
$b = 1$, that is $\beta = \operatorname{vol}
(g)$.
\begin{enumerate}
\item Round $S^2$. Here $c = +\infty$,
since $\beta$ is not exact on the universal
cover $S^2$. Magnetic geodesics are small
circles, all closed.
\item Flat $\mathbb{T}^2$. Again $c =
+\infty$: the lift of $\beta$ to
$\mathbb{R}^2$ is exact, but every primitive
is unbounded. Magnetic geodesics lift to
circles in the plane.
\item A hyperbolic surface $\Sigma_h$. Here
$c(g_{\mathbb{H}},1) = 1$. In the disk
model the picture shows circles below $c$,
horocycles at $c$, and curves meeting the
boundary above $c$. Spoken: below $c$ all
orbits are closed and contractible, with
the same period; at $c$ the flow is the
horocycle flow, which is minimal, so it has
no closed orbits.
\end{enumerate}
\end{example}

\medskip\noindent
{\bf The question.} Given $m > 0$, is there
a closed $(g,\beta)$-geodesic with speed
$m$?
\begin{enumerate}
\item For $m > c$, standard variational
methods apply, and one gets analogues of the
existence results for closed Riemannian
geodesics (attributed on the slides to
G. Contreras). Spoken: there is one in
every free homotopy class.
\item For $m = c$ there may be no closed
orbits. Take a compact surface with
Gaussian curvature $-C < 0$ and $\beta =
\lambda \operatorname{vol}(g)$ with
$\lambda > 0$: at the critical level the
flow is a horocycle flow. A similar
construction works on compact Kähler
hyperbolic manifolds.
\item For $m < c$ the standard techniques
fail. This is the regime of the talk.
\end{enumerate}

\medskip\noindent
{\bf Known results below $c$.} As stated on
the slides; the references were not checked.
\begin{enumerate}
\item Taimanov, and
Contreras--Macarini--Paternain: on a closed
surface with $\beta$ exact, for every
$s \in (0,c]$ there is a closed
$(g,\beta)$-geodesic with speed $s$, not
necessarily contractible.
\item Contreras (also Merry, and
Benedetti--Asselle): for almost every
$m < c$ there is a contractible closed
$(g,\beta)$-geodesic with speed $m$.
\item Ginzburg--Gürel: if $\beta$ is
symplectic, then for every sufficiently
small $m$ there is a contractible closed
$(g,\beta)$-geodesic with speed $m$.
\item Assenza: if $\beta$ vanishes nowhere,
then the same holds for every sufficiently
small $m$. See
Theorem~\ref{theorem-assenza-nonvanishing}.
\end{enumerate}
\emph{Contreras' conjecture} asks for a
contractible closed $(g,\beta)$-geodesic for
every speed $m \in (0,c)$. The value $m = c$
is excluded because of the horocycle
example; spoken, the expectation is that the
Kähler hyperbolic examples are the only
exceptions. Spoken, the almost-every result
comes from a Struwe monotonicity argument,
parametrized by the speed, applied to the
min-max class described below.

\begin{theorem}[Assenza]
\label{theorem-assenza-nonvanishing}
Let $(M,g)$ be closed and let $\beta$ vanish
nowhere. Then for every sufficiently small
speed $m$ there is a contractible closed
$(g,\beta)$-geodesic with speed $m$.
\end{theorem}

This is the preprint arXiv:2309.03159
(Assenza, \emph{Magnetic curvature and
existence of a closed magnetic geodesic on
low energy levels}, revised 2024). The idea
is to bound the period of closed magnetic
geodesics using a \emph{magnetic Ricci
curvature} $\operatorname{Ric}^{\beta,m}$,
which the talk did not define. A
Bonnet--Myers argument gives: if
$\operatorname{Ric}^{\beta,m} \ge \rho^{-2}
> 0$ for some $\rho > 0$, then every closed
magnetic geodesic
$\gamma : [0,T] \to M$ with speed $m$
satisfies
$$
T \le \rho\pi
(\operatorname{index}(\gamma) + 1).
$$
Together with an index estimate, this
upgrades ``almost every speed'' to ``every
speed'' on intervals where the magnetic
curvature is positive. If $\beta$ vanishes
nowhere, $\operatorname{Ric}^{\beta,m}$ is
bounded away from zero for all small $m$.
Then the period stays bounded by some $T^*$,
and the length satisfies
$$
\ell(\gamma) = mT \le mT^* \to 0
$$
as $m \to 0$. So closed magnetic geodesics
with speed $m \to 0$ shrink to points.

\medskip\noindent
{\bf New results.} From now on $M$ need not
be compact and $g$ need not be complete.
Write $|\beta_p|$ for the norm of $\beta$ at
$p$ dual to $g$, and
$$
\|\beta\|_\infty = \sup_{p \in M}
|\beta_p|.
$$

\begin{definition}
\label{definition-local-max}
A \emph{local compact maximum} of $\beta$ is
a subset $K \subset M$ such that
\begin{enumerate}
\item $K$ is compact and $|\beta_p|$ takes
the same value $\|\beta\|_K > 0$ at every
$p \in K$, and
\item there is an open neighborhood $V_K$ of
$K$ with $|\beta_q| < \|\beta\|_K$ for every
$q \in V_K \setminus K$.
\end{enumerate}
\end{definition}

For example, a strict local maximum point of
$|\beta|$ is a local compact maximum, and so
is a compact plateau of $|\beta|$ surrounded
by smaller values. The slides draw both.

\begin{theorem}[Assenza--Benedetti--Macarini]
\label{theorem-abm-local-maximum}
Let $(M,g)$ be a Riemannian manifold and
$\beta$ a magnetic form that is not
identically zero. For every local compact
maximum $K$ of $\beta$ there is $m_K > 0$
such that for every $m \in (0,m_K)$ there is
a contractible closed $(g,\beta)$-geodesic
$\gamma$ with speed $m$ and image in $V_K$.
\end{theorem}

\begin{corollary}[Assenza--Benedetti--%
Macarini]
\label{corollary-abm-closed-manifold}
Let $(M,g)$ be closed and $\beta$ not
identically zero. Then for every
sufficiently small speed $m$ there is a
contractible closed $(g,\beta)$-geodesic.
\end{corollary}

\begin{proof}
On a closed manifold $|\beta|$ attains its
maximum $\|\beta\|_\infty > 0$. The set
where it is attained is compact, and the
neighborhood $V = M$ works in
Definition~\ref{definition-local-max},
since $|\beta| < \|\beta\|_\infty$ off this
set. Apply
Theorem~\ref{theorem-abm-local-maximum}.
\end{proof}

Compactness of $K$ is needed: the slides
mention examples where $K$ is a non-compact
maximum and there are no closed
$(g,\beta)$-geodesics for any small speed.
Spoken: these examples have, the speaker
thinks, an incomplete metric, and it is open
whether there are such examples with a
complete metric. The topology of $K$ is
arbitrary; only compactness matters.

\begin{theorem}[Assenza--Benedetti--Macarini]
\label{theorem-abm-theorem-zero}
Let $(M,g)$ be compact and let $\beta$
vanish nowhere. For every $\varepsilon > 0$
there is $m_\varepsilon > 0$ such that for
every $m \in (0,m_\varepsilon)$ there is a
contractible closed $(g,\beta)$-geodesic
$\gamma$ with speed $m$. Moreover
$\ell(\gamma) \le Cm$ for a constant $C$, and
the image of $\gamma$ lies in
$$
M^\beta_\varepsilon
= \{ q \in M :
|\beta_q| \ge \|\beta\|_\infty - \varepsilon
\}.
$$
\end{theorem}

Spoken, in answer to a question: the
constant $C$ depends only on
$\|\beta\|_\infty$, and is about
$2\pi/\|\beta\|_\infty$, the minimal
period.

On the slides this is ``Theorem 0''. The
passage to
Theorem~\ref{theorem-abm-local-maximum} was
sketched in three steps.
\begin{enumerate}
\item From nowhere vanishing to possibly
vanishing forms, by Thom's transversality
theorem in the space of closed 2-forms.
Spoken: deform $\beta$ away from the
maximum into a nowhere vanishing form; near
the maximum the dynamics does not change.
On a surface this is just a cut-off of the
function $b$.
\item From global to local maxima, by
rescaling the metric and the magnetic form
outside a neighborhood of $K$, so that $K$
becomes a global maximum of the new system.
\item From compact to non-compact manifolds,
by gluing along the boundary of a
neighborhood of $K$ to reduce to a compact
manifold. Spoken: one doubles the
neighborhood along its boundary; the
delicate point is extending $\beta$ to a
closed form on the double.
\end{enumerate}

\medskip\noindent
{\bf Variational setup.} Let $\Lambda_0
\subset H^1(S^1,M)$ be the Hilbert manifold
of contractible loops. For $x \in \Lambda_0$
put
$$
\mathcal{E}(x) =
\left( \int_0^1 |\dot x|^2\, dt
\right)^{1/2},
\quad
\mathcal{B}(x) = \int (D_x)^*\beta,
$$
where $D_x : S^1 \times [0,1] \to M$ is a
capping disk of $x$. In general
$\mathcal{B}$ depends on the capping disk.
Let
$$
\Omega = \{ x \in \Lambda_0 :
0 < \mathcal{E}(x) < d_g \}.
$$
For $d_g$ small enough (spoken: for example
less than twice the injectivity radius),
take capping disks inside short loops. Two
such disks glue to a small sphere, which is
contractible, and $\beta$ is closed, so
$\mathcal{B}$ is single-valued on $\Omega$.
For $m > 0$ define
$\mathcal{A}_m : \Omega \to \mathbb{R}$ by
$$
\mathcal{A}_m(x) = m\mathcal{E}(x)
- \mathcal{B}(x),
\quad
T(x) = \frac{\mathcal{E}(x)}{m}.
$$
The basic fact: $x$ is a critical point of
$\mathcal{A}_m$ with $T(x) = T$ if and only
if
$$
\gamma(t) = x(t/T),
\quad
t \in [0,T],
$$
is a contractible closed magnetic geodesic
with speed $m$. So closed magnetic geodesics
with speed $m$ in $\Omega$ are critical
points of $\mathcal{A}_m$.

\medskip\noindent
{\bf Min-max.} For small $m$ the action
$\mathcal{A}_m$ takes negative values on
$\Omega$: reversing the orientation of a
loop changes the sign of $\mathcal{B}$, and
for $m$ small the term $m\mathcal{E}$ does
not compensate. This is the difference with
the Riemannian energy, which is
non-negative. Consider the paths
$$
\Gamma_m = \{ \psi : [0,1] \to \Omega \},
$$
that start at a loop with
$\mathcal{E}(\psi(0)) = \delta_m$ and end at
a loop with
$\mathcal{A}_m(\psi(1)) < 0$; this class is
non-empty. The min-max value is
$$
u(m) = \inf_{\psi \in \Gamma_m}
\max_{s \in [0,1]}
\mathcal{A}_m(\psi(s)).
$$
An isoperimetric-type inequality
$|\mathcal{B}(x)| \le C\mathcal{E}(x)^2$
gives
$$
\mathcal{A}_m(x) \ge m\mathcal{E}(x)
- D\mathcal{E}(x)^2,
$$
and hence $u(m) > 0$ if $\delta_m$ is small.
So the min-max is non-trivial: every path
must cross a ``mountain'' of positive
action.

\medskip\noindent
{\bf The difficulty.} One wants a
Palais--Smale sequence $x_n$ at level
$u(m)$, meaning
$$
\mathcal{A}_m(x_n) \to u(m),
\quad
\|d\mathcal{A}_m(x_n)\| \to 0,
$$
which is compact. Compactness needs a bound
on $\mathcal{E}(x_n)$, that is on the length,
and nothing in $\mathcal{A}_m$ alone gives
it. Struwe's monotonicity argument, which
gives such control for almost every $m$,
does not apply here.

\medskip\noindent
{\bf Controlling $\mathcal{E}$.} For
$m^* > m$,
$$
\mathcal{A}_{m^*}(x) - \mathcal{A}_m(x)
= (m^* - m)\mathcal{E}(x),
$$
so controlling $\mathcal{A}_m$ and
$\mathcal{A}_{m^*}$ together controls
$\mathcal{E}$. The plan is a vector field
$W_m$ that is a pseudogradient for both
$\mathcal{A}_m$ and $\mathcal{A}_{m^*}$, for
example by interpolating between their
gradients. This fails exactly where
$d\mathcal{A}_m$ and $d\mathcal{A}_{m^*}$
vanish or point in opposite directions,
that is where
$$
\lambda\, d\mathcal{A}_{m^*}(x)
+ (1 - \lambda)\, d\mathcal{A}_m(x) = 0
$$
for some $\lambda \in [0,1]$. Since
$\mathcal{A}_m$ is affine in $m$, the left
side is $d\mathcal{A}_\mu(x)$ with
$$
\mu = \lambda m^* + (1 - \lambda)m,
$$
so $x$ is a critical point of
$\mathcal{A}_\mu$ for an intermediate speed
$\mu \in [m,m^*]$.

\medskip\noindent
{\bf Bonnet--Myers again.} Fix
$\mathcal{E}_0 \in (0,d_g)$. If $m$ and
$m^*$ are small enough, every critical point
$x$ of $\mathcal{A}_\mu$ with
$\mu \in [m,m^*]$ and
$\mathcal{E}(x) \ge \mathcal{E}_0$ has
$$
\operatorname{index}_\mu(x) \ge 3.
$$
The idea is to push min-max paths, which
are 1-dimensional, off these critical points
using negative directions of the Hessian of
$\mathcal{A}_\mu$. Spoken, why index 3: by
the difference formula above, both actions
decrease together along deformations that
keep $\mathcal{E}$ constant. There are no
arbitrarily short closed geodesics, so
$d\mathcal{E} \neq 0$ there, and by the
implicit function theorem the level set of
$\mathcal{E}$ is a hypersurface. A negative
subspace of dimension 3 still meets it in
dimension at least 2, which is enough to
push 1-dimensional paths off. The final
pseudogradient for $\mathcal{A}_m$ on
$\Gamma_m$ interpolates between $W_m$ and
the gradient of $\mathcal{A}_m$.

\medskip\noindent
{\bf Confinement on a surface.} Let $M$ be a
surface and $b$ a nowhere vanishing magnetic
function. Then
$$
\mathcal{A}_m(x) = m\mathcal{E}(x)
- \int_{S^1 \times [0,1]}
(D_x)^* b \operatorname{vol}(g).
$$
\begin{enumerate}
\item Upper bound. Evaluating
$\mathcal{A}_m$ on small concentric circles
around a point where $|b|$ is maximal gives
$$
u(m) \le
\frac{\pi m^2}{\|b\|_\infty}
+ o(m^2)
$$
as $m \to 0$.
\item Lower bound. Let $x_m$ be the critical
points of $\mathcal{A}_m$ at level $u(m)$
from the min-max. By the Bonnet--Myers bound
they shrink to a point $p \in M$. Spoken:
the periods are bounded above and below, so
the unit velocities
$\zeta_m = \dot\gamma_m/m$, curves in the
unit sphere bundle, converge uniformly to a
solution of the equation in the fiber at
$p$, where the Lorentz force is
$|b(p)|J$. Its solutions are harmonic
oscillators with period $2\pi/|b(p)|$. In
normal coordinates at $p$ the velocity is,
to first order,
$\zeta(t) = z_0 e^{i|b(p)|t}$, and
integrating gives
$$
\gamma_m(t) =
-m\frac{e^{i|b(p)|t}}{|b(p)|} + o(m),
\quad
t \in \left[0,
\frac{2\pi}{|b(p)|}\right].
$$
This is a circle of radius $m/|b(p)|$ and
period $T_m = 2\pi/|b(p)| + o(1)$. Using
$\mathcal{E}(x_m) = mT_m$,
$$
\begin{aligned}
\mathcal{A}_m(x_m)
&= m^2T_m - |b(p)|
\frac{\pi m^2}{|b(p)|^2} + o(m^2) \\
&= \frac{2\pi m^2}{|b(p)|}
- \frac{\pi m^2}{|b(p)|} + o(m^2) \\
&= \frac{\pi m^2}{|b(p)|} + o(m^2).
\end{aligned}
$$
\end{enumerate}
Since $\mathcal{A}_m(x_m) = u(m)$, the two
estimates give
$$
\frac{\pi m^2}{|b(p)|}
\le \frac{\pi m^2}{\|b\|_\infty}
+ o(m^2),
$$
hence $|b(p)| \ge \|b\|_\infty$, so
$|b(p)| = \|b\|_\infty$. The orbits
concentrate at the maximum of $|b|$, which
is the localization in
Theorem~\ref{theorem-abm-theorem-zero}.

\begin{remark}
\label{remark-magnetic-slide-corrections}
Changes from the slides: the slides call the
action $\mathcal{A}_k$ in places, write the
interpolation as
$\lambda\, d\mathcal{A}_{m^*}
+ (1-\lambda)\mathcal{A}_m$ without the
second $d$, and write the interval as
$[m^*,m]$ although $m^* > m$. The last
slide writes the chain of equalities and the
inequality in one line; the inequality
direction above is my reading of how the
two ingredients combine.
\end{remark}

\medskip\noindent
{\bf Questions (spoken).}
\begin{enumerate}
\item Long closed orbits exist, very slow,
near the zero set of $|b|$, where the
magnetic action is weak. Away from it, as
$m \to 0$ orbits shrink to points.
\item Each strict local maximum of $|\beta|$
gives a closed orbit for small speed, so
their number gives a lower bound; it
depends on $\beta$, and on a surface $b$ can
have a single maximum.
\item Work in progress: multiplicity results
via Lusternik--Schnirelmann theory, using
min-max over homology classes of $M$
instead of local maxima. In the nowhere
vanishing case the expectation is at least a
couple of closed orbits.
\end{enumerate}

\bibliography{my}
\bibliographystyle{amsalpha}

\end{document}
